\chapter{Vector bundles and complex manifolds} \label{appendix_Kahler}
This chapter presents a brief introduction to the theory of complex and Kähler manifolds, both of them will constitute the basic frame in which we will work. We also devote part of the chapter to introducing abstract vector bundles. In basic differential geometry, one usually learns about the tangent and cotangent bundles, as well as tensor bundles. These ideas can be generalized to abstract vector bundles, which constitute the basic framework we will need for chapter \ref{ch_connections}. Since this is understood as preliminary material, we will omit most of the proofs. We refer the reader to the books \parencite{lee_introCM} and \parencite{huybrechts2005complex}, for everything related to complex and Kähler manifolds. For the basics of vector bundles, \parencite[Ch. 10]{lee_introSM} is a good reference. It only deals with real smooth vector bundles, but most things apply to complex and holomorphic vector bundles with minor modifications. Since we will use \(i\) as an index for coordinate expressions very often, we set the notation \(\ii:=\sqrt{-1}\). We introduce in this section the Dolbeault complex, i.e. complex differential forms. To differentiate them from real differential forms, let us set the notation \(\Omega_\R^p(M)\) for the \(\mathscr{C}^\infty(M;\R)\)-module of \(p\)-differential forms on a smooth manifold \(M\).

\section{Complex manifolds}
Let us start by recalling some facts from complex analysis, and introducing some concepts from multivariable complex calculus. Recall that if \(f\colon W \to \C\) is a function, where \(W \subset \C\) is open, it is \textit{holomorphic} if it has a complex derivative at any point of \(W\), or, equivalently, if, when writing \(f=u+\ii v\), as seeing it as a map from \(\R^2\) to \(\R^2\), its real and imaginary parts satisfy the \textit{Cauchy-Riemann} equations, i.e \(\partial_x u=\partial_y v\) and \(\partial_y u=-\partial _x v\). Equivalently, if we introduce the operators
\begin{equation} \label{derparz}
\derpar{}{z}:=\frac{1}{2} \parent{\derpar{}{x}-\ii \derpar{}{y}} \hspace{1 cm}
\derpar{}{\bar{z}}:=\frac{1}{2}\parent{\derpar{}{x}+\ii \derpar{}{y}},
\end{equation}
then \(f\) is holomorphic if and only if \(\partial_{\bar{z}}f=0\) and \textit{antiholomorphic} (that is, \(\bar{f}\) is holomorphic) if and only if \(\partial_zf=0\). These definitions can be adapted to functions of several complex variables: if \(W \subset \C^n\) is open and \(f\colon W \to \C^m\) is continuous in \(W\) and, for each \(p \in W\), the following limit exists
\begin{equation} \label{lim_derpar}
\derpar{f}{z^j}:=\lim_{h \to 0} \frac{f(p^1, \ldots , p^j+h, \ldots p^n)-f(p^1, \ldots p^n)}{h}
\end{equation}
then we say that \(f\) is holomorphic in \(W\). The Cauchy-Riemann equations still give a characterization of holomorphicity (but now they have to be satisfied for all coordinates). Extending the definition of \eqref{derparz} to several complex variables, we have 
\[
\derpar{f}{z^j}=\derpar{f}{x^j}=\frac{1}{\ii} \cdot \derpar{f}{y^j},
\]
which coincides with the limit in equation \eqref{lim_derpar} when calculated at a point where \(f\) is holomorphic. We will use, when needed, the identification \(\R^{2n} \equiv \C^n\) via \((x^1, y^1, \ldots , x^n, y^n) \leftrightarrow (x^1+\ii y^1, \ldots x^n+\ii y^n)\), and we set \(z^j:=x^j + \ii y^j\). \\

Now, the definition of a complex manifold is the obvious modification of the definition of a smooth manifold:

\Definition{}{A \textit{complex manifold} \(M\) is a second countable Hausdorff topological space, endowed with a complex atlas, that is, a family \(\mathscr{A}\) of pairs (complex charts) \((U, \varphi)\) where each \(U \subset M\) is open, the open sets cover \(M\), \(\varphi \colon U \to \C^n\) is a homeomorphism, and whenever \((U, \varphi), (V, \psi) \in \mathscr{A}\) satisfy, \(U \cap V \neq \emptyset\) the maps,
\[
\psi \circ \varphi^{-1}\colon \varphi(U \cap V) \to \psi(U \cap V) \hspace{1 cm}
\varphi \circ \psi^{-1} \colon \psi(U \cap V) \to \varphi(U \cap V) 
\]
are holomorphic.
}
Then, a map between complex manifolds \(f\colon M \to N\) is holomorphic, if its \textit{localization}\index{Localization} is holomorphic, that is, for all charts \((U, \varphi)\) and \((V, \psi)\) of \(M\) and \(N\) respectively, the map \(\psi \circ f \circ \varphi^{-1}\) is holomorphic. A holomorphic bijective map whose inverse is also holomorphic is called a \textit{biholomorphism}\index{Biholomorphism}. Whenever \(M\) is a complex manifold, \(\mathcal{O}_U\) denotes the ring of all holomorphic functions on the open set \(U\). Observe that all complex manifolds are even-dimensional when considered as real manifolds. Moreover, one may show that complex manifolds are always orientable (see \parencite[Prop. 1.49]{lee_introCM} for a proof). \\

The concepts of \textit{submersion, immersion} and \textit{embedding} admit direct generalizations to complex manifolds: a holomorphic map is said to have either of these properties if it has them when considered as a smooth map between the smooth underlying manifolds.

\begin{example}
    Some examples of complex manifolds are \(\C^n\) and its open subset and \(\CP^n\). More generally, every open subset of a complex manifold inherits the structure of a complex manifold.
\end{example}
In an analogy with smooth manifolds, we will often say ``let \((z_1, \ldots z_n)\) be local holomorphic coordinates'' of a complex manifold \(M\), meaning we have chosen a complex chart \((U, \varphi)\), where the components of \(\varphi\) are called \(z^j\).

\begin{obs}
    The notion of Lie group admits a direct generalization to complex manifolds, which is as follows: A \textit{complex Lie group}\index{Lie group!complex} is a complex manifold \(G\) such that the group structure is compatible with the complex structure, i.e the maps \(G\times G \to G, (g,h) \mapsto g \cdot h^{-1}\) and \(G \to G, g \mapsto g^{-1}\) are holomorphic. For example, \(\mathrm{GL}(n,\C)\) is a complex manifold (it can be seen as the open set \(\det^{-1}(\C^*) \subset \C^{n^2}\)), and the entries of the product and inverse are either polynomial or rational maps, both of which are holomorphic. Of course, in this case, Lie group homomorphisms are required to be \textit{holomorphic} group homomorphism (not merely smooth).
\end{obs}
\section{Vector bundles}
% We devote this chapter to the study of vector bundles. The exposition of the section is, with some minor changes to include holomorphic vector bundles, chapter 10 of \parencite{lee_introSM}. This can be done form a purely topological viewpoint, but since we will be using holomorphic vector bundles very frequently throughout the text, we shall develop the theory for the latter. Nonetheless, the theory of holomorphic vector bundles is not that different from the theory of smooth vector bundles. Aside from the occasional swap of the word "holomorphic" by the word "smooth". This means that both theories can be developed simultaneously. 
We will use both smooth and holomorphic vector bundles, but given how similar these concepts are, we will introduce them simultanously. To do so, we shall mimic the notation used in \parencite{garcia2007differential}. We will use the symbol \(\mathcal{S}\) to denote either smooth or holomorphic functions, sections, maps etc. For example, \(\mathcal{S}(M)\) denotes either the set of smooth functions on a manifold, \(\mathscr{C}^\infty(M;\K)\) (\(\K=\R\) or \(\C\), depending on the context), or the set of holomorphic functions on a complex manifold, \(\mathcal{O}_M\). Similarly, an \(\mathcal{S}\)-manifold is either a smooth manifold, or a holomorphic one, an \(\mathcal{S}\)-morphism between \(\mathcal{S}\)-manifolds, is either a smooth map, or a holomorphic one, and an \(\mathcal{S}\)-isomorphism is either a diffeomorphism or a biholomorphism. Finally, an \(\mathcal{S}\)-manifold is either a smooth or complex manifold. \\

With these ideas in mind, we are ready to define the notion of vector bundle. Vector bundles appear as a generalization of the tangent bundle of a smooth manifold. Intuitively, the tangent bundle of a manifold is obtained by gluing the tangent space to the manifold at a point, \(T_pM\) to the corresponding point, \(p\), and then considering the union of all those spaces. Vector bundles try to generalize this concept by replacing the tangent spaces \(T_pM\) by arbitrary vector spaces \(E_p\) (all of them of the same dimension). Therefore, now we have to glue a vector space \(E_p\) to each point \(p\), and consider the union of all those spaces. Of course, the added generality means that now there is not a canonical topology or smooth structure (in general), so those should be given in advance. \\

% This is very useful because it allows us to consider all tangent spaces as a single geometric object, and when given its natural topology and smooth structure, it makes sense of the intuitive notion of smooth vector field; that is, an assignment of a vector to each point of the manifold which varies smoothly with the point. \\
% Now that we have in mind what we want to define, let us proceed with the formal definitions.

\Definition{\textbf{Vector bundle}}{Let \(M\) be an \(\mathcal{S}\)-manifold. A real (or complex) \(\mathcal{S}\)-vector bundle of rank \(k\) over \(M\) is another \(\mathcal{S}\)-manifold, \(E\), together with a surjective \(\mathcal{S}\)-morphism \(\pi\colon E \to M\) satisfying the following conditions
\begin{enumerate}
    \item For each \(p \in M\), \(E_p:=\pi^{-1}(p)\) is a \(k\)-dimensional real (or complex) vector space.
    \item For each \(p \in M\) there is an open neighborhood \(U\) of \(p\), and an \(\mathcal{S}\)-isomorphism \(\Phi\colon \pi^{-1}(U)\to U \times \R^k\) (or \(\Phi\colon \pi^{-1}(U)\to U \times \C^k\)), called a \textit{local trivialization} satisfying the following conditions:
    \begin{enumerate}
        \item \(\pi_U \circ \Phi=\pi\), where \(\pi_U \colon U \times \R^k \to U\) is the usual projection;
        \item \(\Phi|_{E_q}\colon E_q \to \corch{q} \times \R^k\) (or \(\Phi|_{E_q}\colon E_q \to \corch{q} \times \C^k\)) is a linear isomorphism.
    \end{enumerate} \index{Vector bundle} \index{Local!trivialization}
\end{enumerate}
\label{def_vector_bundle}}
The space \(M\) is often called the \textit{base space}, and \(E\) is called the \textit{total space}. The map \(\pi\colon E \to M\) is called the \textit{projection}. From now on, the symbol \(\K\) will be used to denote either \(\R\) or \(\C\). Let us make clear that the field may still be \(\C\) even when we consider smooth not holomorphic bundles. Indeed, there are a lot of interesting complex smooth vector bundles that do not admit a holomorphic structure such as \(T_\C^*M\) (see section \ref{section_complexified_bundles} below). Finally, observe that we do not restrict ourselves to consider only holomorphic vector bundles over complex manifolds, sometimes we need to study complex smooth bundles over those as well. The dimension of each fiber on a vector bundle is called its \textit{rank}\index{Rank of a vector bundle}, and it is denoted \(\mathrm{rank}(E)\). \\

\begin{example} \label{ex_trivial_bundle}
    Let \(M\) be an \(\mathcal{S}\)-manifold. Then, \(M \times \K^r\) is a \(\K\)-vector bundle over \(M\). The projection is the trivial one, and it has one single trivialization, \(\Phi=\mathrm{id}_{M \times \K^r}\).
\end{example}

The previous example is a particular case of what we call a \textit{trivial bundle}.\index{Vector bundle!trivial} A trivial bundle is a bundle for which there is a global trivialization, that is, a trivialization defined over all of \(M\). This of course implies that if a bundle \(E\) over \(M\) is trivial, then it is \(\mathcal{S}\)-isomorphic to \(M \times \R^k\) (or \(\C^k\) if the bundle was complex). Once we define the right notion of morphism between vector bundles, we will see that this in fact implies that \(E\) and \(M \times \R^k\) are isomorphic.  \\

% The meaning of the existence of local trivializations is that, even though the bundle will not be trivial in general, it is union of open sets that are trivial. In that sense, trivializations fulfill a similar role as charts do in manifolds, and as such, there is a change of trivialization formula, as shown by the following lemma.

% \Lemma{}{Let \(\pi\colon E \to M\) be a real or complex range \(k\) \(\mathcal{S}\)-vector bundle. Let \((U, \Phi)\) and \((\Psi, V)\) two local trivializations such that \(U \cap V \neq \emptyset\). Then, there is an \(\mathcal{S}\)-morphism \(\tau\colon U \cap V \to \mathrm{GL}(k, \K)\) such that
% \[
% \parent{\Phi \circ \Psi^{-1}}(p, v) = \parent{p, \tau(p)v} \hspace{1 cm} p \in U \cap V, \ \ v \in \K^r,
% \]
% where \(\tau(p)v\) is the product of the matrix \(\tau(p)\) and the vector \(v \in \K^r\). \label{lemma_existence_tau}
% }
% \begin{proof}
%     See \parencite[Lemma 10.5]{lee_introSM}.
% \end{proof}

The maps \(\tau\colon U \cap V \to \mathrm{GL}(k, \K)\) are called \textit{transition functions}\index{Transition functions}, even though, strictly speaking, they are not functions, but maps whose codomain is \(\mathrm{GL}(n, \K)\). \\

It is often the case that we want to create a vector bundle by gluing a family of vector spaces we choose to each point of a manifold. However, this alone does not provide the union with the structure of an \(\mathcal{S}\)-vector bundle, as we still have to define a topology and an \(\mathcal{S}\)-structure on the union. If we were to this directly every time, we would have to define the topology, the \(\mathcal{S}\)-structure and the trivializations, and verify all the requirements of the definition of vector bundle. The following lemma provides a shortcut; if we fix the trivializations, and they have transition functions that are \(\mathcal{S}\)-morphisms, then both the topology and the \(\mathcal{S}\)-structure are determined unequivocally. \\

Before stating the theorem, let us briefly recall that if we are given an index set \(A\), and a family of sets indexed by \(A\), \(\corch{W_a}_{a \in A}\), the disjoint union is defined as
\begin{equation} \label{eq_disjoint_union}
\bigsqcup_{a \in A} W_a:=\bigcup_{a \in A} \corch{a} \times W_a.
\end{equation}
In particular, elements of the disjoint union are pairs \((a, w)\) such that \(w \in W_a\). Under some abuse of notation, it is quite common to omit the indexing element, and not write elements of the disjoint union.

\Lemma{\textbf{Vector bundle chart lemma}}{Let \(M\) be an \(\mathcal{S}\)-manifold, and suppose we are given a vector space \(E_p\) for each \(p \in M\), all of them of the same dimension \(r\). Define \(E:=\bigsqcup_{p \in M}E_p\), and let \(\pi \colon E \to M, (p,v) \mapsto v\). Suppose as well that we are given the following data: \label{vector_bund_chart_lemma}
\begin{enumerate}
    \item an open cover \(\corch{U_\alpha}_{\alpha \in A}\) of \(M\),
    \item for each \(\alpha \in A\), a bijection \(\Phi_\alpha\colon \pi^{-1}(U_\alpha) \to U_\alpha \times \K^r\) that restricts to a vector space isomorphism \(\Phi_\alpha|_{E_p}\colon E_p \to \corch{p} \times \K^r\) for each \(p \in U_\alpha\),
    \item for each pair \(\alpha, \beta \in A\) such that \(U_\alpha \cap U_\beta \neq \emptyset\), an \(\mathcal{S}\)-morphism, \(\tau_{\alpha\beta}\colon U_\alpha \cap U_\beta \to \mathrm{GL}(k, \K)\), such that
    \begin{equation} \label{eq_def_tau}
    \begin{array}{cccc}
        \Phi_\alpha \circ \Phi_\beta^{-1} \colon & \parent{U_\alpha \cap U_\beta} \times \K^r & \longrightarrow & \parent{U_\alpha \cap U_\beta} \times \K^r \\
        & (p,v) & \mapsto & (p, \tau(p)v).
    \end{array}
    \end{equation}
    Then, \(E\) has a unique topology and smooth structure making it into an \(\mathcal{S}\)-manifold, and a rank \(r\) \(\mathcal{S}\)-vector bundle over \(M\) such that \((U_\alpha, \Phi_\alpha)\) is a trivialization for every \(\alpha \in A\).
\end{enumerate}
}
\begin{proof}
    See \parencite[Lemma 10.6]{lee_introSM}
\end{proof}

\begin{example} \label{ex_vector_bundle_constructions}
    Lemma \ref{vector_bund_chart_lemma} is very useful because it allows to create a wide range of vector bundles out of existing ones. More precisely, it can be used to "elevate" a lot of constructions that can be done with vector spaces, such as sums, tensor products etc, to bundles. We will present the explicit construction of the tensor product bundle. All other bundles in this example can be constructed by similar procedures.
    \begin{enumerate}
        \item Let us show how the \textit{tensor product bundle}\index{Vector bundle!tensor product} is defined. Let \(E, F\) be two \(\mathcal{S}\)-vector bundles over the same base space \(M\), of respective ranks \(r\) and \(s\). We want to define another vector bundle whose at each point is \(E_p \otimes F_p\). To do so, we can use lemma \ref{vector_bund_chart_lemma} as follows. Define \(E \otimes F:=\bigsqcup_{p \in M} E_p \otimes F_p\) and the obvious projection. For the open cover, consider all intersections \(U \cap V\) such that \((U, \Phi)\), \((V, \Psi)\) are local trivializations of \(E \to M\) and \(F \to M\) respectively. Define trivializations \(\Gamma_{\Phi,\Psi}\) for \(E \otimes F\) by linear extensions in the fibers of
    \[
    \begin{array}{cccc}
          \pi^{-1}(U \cap V) & \longrightarrow & (U \cap V) \times \K^{r\cdot s}\\
         (p, v, w) & \mapsto &(p, \pi_r(\Phi(p,v)), \pi_s(\Psi(p,v))),
    \end{array}
    \] 
    where \(\pi_r\colon U \times \K^r \to \K^r\) is the usual projection. These are bijections because \(\K^{r\cdot s} \cong \K^r \otimes \K^s \cong E_p \otimes F_p\), and \(\Gamma_{\Phi,\Psi}\) restricts to an isomorphism on each fiber by definition. Only the third condition of lemma \ref{vector_bund_chart_lemma} needs to be verified. Since \(E\) and \(F\) are \(\mathcal{S}\)-vector bundles over \(M\), if we are given \((U_1, \Phi_1)\) and \((U_2, \Phi_2)\) two local trivializations such that \(U_1 \cap U_2 \neq \emptyset\), then there is an \(\mathcal{S}\)-morphism \(\tilde{\tau}\colon U_1 \cap U_2 \to \mathrm{GL}(k, \K^r)\) satisfying the condition of equation \eqref{eq_def_tau} (see lemma \parencite[Lemma 10.5]{intro_leeSM}), that is,
    \[
    \parent{\Phi_1 \circ \Phi_2^{-1}}(p, v) = (p, \tau(p)v)
    \]
    and the same applies for two local trivializations \((V_1, \Psi_1)\) and \((V_2, \Psi_2)\) such that \(V_1 \cap V_2 \neq \emptyset\) and a map \(\hat{\tau}\). Then, if \((U_1 \cap U_2) \cap (V_1 \cap V_2) \neq \emptyset\), we have
    \[
    \Gamma_{\Phi_1 \Psi_1} \circ \Gamma_{\Phi_2 \Psi_2}^{-1}(p,u) = (p, \left[\tau(p) \otimes \hat{\tau}(p)\right](u)),
    \]
    where \(\tau(p) \otimes \hat{\tau}(p)\) is the \textit{tensor product of linear maps}, i.e. the linear extension of the map acting on decomposable elements as \(\left[\tau(p) \otimes \hat{\tau}(p)\right] (v \otimes w)= \left[\tau(p)v \right] \otimes \left[\hat{\tau}(p) w\right]\) (this map is well-defined by the universal property of tensor products). Indeed,
    \begin{align*}
        \Gamma_{\Phi_1 \Psi_1} \circ \Gamma_{\Phi_2 \Psi_2}^{-1}(p, v \otimes w) & = \Gamma_{\Phi_1 \Psi_2}(p, \Phi_2^{-1}(p,v)\otimes \Psi_2^{-1}(p,v)) \\
        & = \Gamma_{\Phi_1 \Psi_1}(p, \pi_r(\Phi_1\circ \Phi_2^{-1}(p,v)) \otimes \pi_2(\Psi_1 \circ \Psi_2(p,w))) \\
        & = (p, \left[\tau(p)v \right] \otimes \left[\hat{\tau}(p) w\right])
    \end{align*}
    and by linear extension, we see that this works for every vector in \(\K^s \otimes \K^r\). Finally, \(\tau \otimes \hat{\tau}\colon U \cap V \to \mathrm{GL}(r, \K) \otimes \mathrm{GL}(s, \K) \cong \mathrm{GL}(r\cdot s, \K)\), and it is an \(\mathcal{S}\)-morphism (this is essentially due to the way the tensor product of matrices is defined see \parencite[Section 11.2]{dummit2003abstract}).
    \item Given a vector bundle \(E\), its \textit{dual vector bundle}\index{Vector bundle!dual}, \(E^*\) is the vector bundle whose fiber at each point \(p \in M\) is \(E_p^*\). Again, one uses lemma \ref{vector_bund_chart_lemma} and checks that the transition functions for this bundle are the transpose inverse of those of \(E\).
    \item Given two vector bundles, \(E\) and \(F\), \(\mathrm{Hom}(E,F)\) is the bundle whose fiber at each point is the vector space of linear homomorphisms \(\mathrm{Hom}(E_p, F_p)\). When \(E=F\), we denote \(\mathrm{Hom}(E,E)=\operatorname{End}(E)\).
        % The fact that \(\tilde{\tau}\) and \(\hat{\tau}\) are \(\mathcal{S}\)-morphisms implies that so is \(\tau\), and using lemma \ref{vector_bund_chart_lemma}, we conclude that \(E \oplus F\) has a unique \(\mathcal{S}\)-vector bundle structure making the maps \(\Gamma_{\Phi, \Psi}\) intro local trivializations.

        \item Given two vector bundles \(E\) and \(F\) one may construct its \textit{Whitney sum}\index{Vector bundle!Whintey sum}, i.e the vector bundle \(E \oplus F\) whose fiber at each point \(p \in M\) is \(E_p \oplus F_p\). The construction follows the same ideas as the tensor product. Nonetheless, it is done explicitly in \parencite[Example 10.7]{lee_introSM}
    % \item Consider for example the sum (also called \textit{Whiney sum}\index{Whintey sum}). Let \(E, F\) be two \(\mathcal{S}\)-vector bundles over the same base space \(M\), of respective ranks \(r\) and \(s\). We want to define another vector bundle whose at each point is \(E_p \oplus F_p\). To do so, we can use lemma \ref{vector_bund_chart_lemma} as follows. Define \(E \oplus F:=\bigsqcup_{p \in M} E_p \oplus F_p\) and the obvious projection. For the open cover, consider all intersections \(U \cap V\) such that \((U, \Phi)\), \((V, \Psi)\) are local trivializations of \(E \to M\) and \(F \to M\) respectively. Define trivializations for \(E \oplus F\) as follows
    % \[
    % \begin{array}{cccc}
    %     \Gamma_{\Phi,\Psi}\colon & \pi^{-1}(U \cap V) & \longrightarrow & (U \times V) \times \K^{r+s}\\
    %     & (p, v, w) & \mapsto &(p, \pi_r(\Phi(p,v)), \pi_s(\Psi(p,v))),
    % \end{array}
    % \] 
    % where \(\pi_r\colon U \times \K^r \to \K^r\) is the usual projection. These maps are bijections because \(\K^{r+s} \cong \K^r \oplus \K^s\), and they clearly restrict to linear isomorphisms from \(E_p \oplus F_p\) to \(\corch{p} \times \K^{r+s}\). Only the third condition of lemma \ref{vector_bund_chart_lemma} needs to be verified. Since \(E\) and \(F\) are \(\mathcal{S}\)-vector bundles over \(M\), if we are given \((U_1, \Phi_1)\) and \((U_2, \Phi_2)\) two local trivializations such that \(U_1 \cap U_2 \neq \emptyset\), then there is an \(\mathcal{S}\)-morphism \(\tilde{\tau}\colon U_1 \cap U_2 \to \mathrm{GL}(k, \K^r)\) satisfying the condition of lemma \ref{lemma_existence_tau}, that is,
    % \[
    % \parent{\Phi_1 \circ \Phi_2^{-1}}(p, v) = (p, \tau(p)v)
    % \]
    % and the same applies for two local trivializations \((V_1, \Psi_1)\) and \((V_2, \Psi_2)\) such that \(V_1 \cap V_2 \neq \emptyset\) and a map \(\hat{\tau}\). Then, if \((U_1 \cap U_2) \cap (V_1 \cap V_2) \neq \emptyset\), we have
    % \[
    % \Gamma_{\Phi_1 \Psi_1} \circ \Gamma_{\Phi_2 \Psi_2}^{-1}(p,v, w) = (p, \tau(p)(v,w))
    % \]
    % where, 
    % \[
    % \tau(p)=\begin{pmatrix}
    %     \tilde{\tau}(p) & 0 \\
    %     0 & \hat{\tau}(p)
    % \end{pmatrix}.
    % \]
    % Indeed, 
    % \begin{equation} \nonumber
    %     \begin{split}
    %         \Gamma_{\Phi_1 \Psi_1} \circ \Gamma_{\Phi_2 \Psi_2}^{-1}(p,v, w) & =  \Gamma_{\Phi_1 \Psi_1} \parent{p, \Phi_2^{-1}(p,v), \Psi_2^{-1}(p,v)} = (p, \pi_r \big(\parent{\Phi_1 \circ \Phi_2^{-1}}(p,v) \big), \pi_s\big(\parent{\Psi_1 \circ \Psi_2^{-1}}(p,v) \big) ) \\
    %         & = (p, \tilde{p}v, \hat{\tau}(p) w).
    %     \end{split}        
    % \end{equation}
    % The fact that \(\tilde{\tau}\) and \(\hat{\tau}\) are \(\mathcal{S}\)-morphisms implies that so is \(\tau\), and using lemma \ref{vector_bund_chart_lemma}, we conclude that \(E \oplus F\) has a unique \(\mathcal{S}\)-vector bundle structure making the maps \(\Gamma_{\Phi, \Psi}\) intro local trivializations.
    % \item The restriction of a vector bundle to a submanifold is again a vector bundle. \index{Vector bundle!restriction of}Let \(\pi\colon E \to M\) be a rank \(r\) \(\mathcal{S}\)-vector bundle over an \(\mathcal{S}\)-manifold \(M\), and let \(S \subset M\) be an immersed or embedded \(\mathcal{S}\)-submanifold. We define the \textit{restriction of \(E\) to \(S\)} as the set \(E|_S:=\bigsqcup_{p \in S} E_p\) and set the projection as the restriction to \(\pi^{-1}(S)\) of the original projection,  \(\pi_S:=\pi|_{\pi^{-1}(S)}\). Let \(\corch{(U_\alpha, \Phi_\alpha)}_{\alpha \in A}\) be the set of trivializations of \(E\). As \(S\) is a submanifold, \(\corch{U_\alpha \cap S}_{\alpha \in A}\) is an open cover of \(S\) (this is obvious for the embedded case, but it remains true for the immersed one, \parencite[Ex 5.20]{lee_introSM}). Define local trivializations by simply restricting local trivializations of \(E\). That is, for a local trivialization \((U_\alpha, \Phi_\alpha)\) define
    % \[
    % \hat{\Phi}_\alpha:=\Phi_\alpha|_{\pi^{-1}(U_\alpha \cap S)}\colon \pi^{-1}(U_\alpha \cap S) \longrightarrow (U_\alpha \cap S) \times \K^r.
    % \]
    % For each \(p \in S\), \(\hat{\Phi}_\alpha|_{E_p}=\Phi_\alpha|_{E_p}\), so the restriction is an isomorphism, and the composition of two such trivializations is
    % \[
    % \parent{\hat{\Phi}_\beta \circ \hat{\Phi}_\alpha^{-1}}(p,v) = (\Phi_\beta \circ \Phi_\alpha^{-1})(p,v) \hspace{1 cm} \text{for} \ \ p \in U_\alpha \cap U_\beta \cap S.
    % \]
    % Once again, the hypothesis of lemma \ref{vector_bund_chart_lemma} are fulfilled, and \(E|_S\) is an \(\mathcal{S}\)-vector bundle over \(S\). 
    \item Let \(E\) be a complex vector bundle, we may define its \textit{complex conjugate vector bundle}\index{Vector bundle!complex conjugate}, \(\bar{E}\) as follows. By definition, the fibers of \(\bar{E}\) are the complex conjugate vector spaces to the fibers of \(E\). More precisely, for each \(x \in E\), \(\bar{E}_x\) coincides with \(E_x\) as a real vector space, but we see it as a complex vector space by defining its product by a complex number as follows 
    \[
    \lambda \ *_{\bar{E}} v:\overline{\lambda} \cdot v,
    \]
    where the second product is the usual product by a complex number in \(E_x\). This turns \(\bar{E}_x\) into a complex vector space, and the set \(\sqcup_{x \in M} \corch{x} \times \bar{E}_x\) with the obvious projection into a smooth complex vector bundle over \(M\).
    \item The quotient of two \(\mathcal{S}\)-bundles in the bundle whose fiber at each point is the quotient of fibers.
    \end{enumerate}
\end{example}

There is a natural notion of subbundle.

\Definition{\textbf{Subbundle}}{Let \(\pi\colon E \to M\) be an \(\mathcal{S}\)-bundle over an \(\mathcal{S}\)-manifold \(M\). An \(\mathcal{S}\)-subbundle is an embedded submanifold \(D \subset E\) such that \(D_p:=D \cap E_p\) is a vector subspace of \(E_p\) for each \(p \in M\), all of them of the same dimension, together with the restriction \(\pi|_D\colon D \to M\), such that \(\pi|_D\colon D \to M\) is a vector bundle.}

Finally, there is natural notion of pullback as well. Consider two \(\mathcal{S}\)-manifolds \(M\) and \(N\), with an \(\mathcal{S}\)-morphism \(\varphi\colon M \to N\), and let \(E\to N\) be an \(\mathcal{S}\)-vector bundle. Then, we may define the \textit{pullback vector bundle}
\begin{equation} \label{def_pullback_vec_bund}
    \varphi^*E:=\corch{(m, v)\in M \times E: \varphi(m)=\pi(v)}.
\end{equation}
This has a natural vector bundle structure with the projection given by \((m,v)\mapsto m\). The fiber over each point is \(\corch{m}\times E_{\varphi(m)}\), and applying the vector bundle chart lemma (lemma \ref{vector_bund_chart_lemma}) we may see that it is an \(\mathcal{S}\)-vector bundle over \(M\). See \parencite[Page 75]{lee_introCM} for details.

\subsection{Local and global sections}
In the usual tangent bundle, vector fields are thought of as maps sending each point to a vector and varying smoothly. More formally, they are sections of \(TM\). The analogous concept for vector bundles is that of a section.

\Definition{\textbf{Section}}{Let \(\pi\colon E \to M\)be an \(\mathcal{S}\)-bundle. A \textit{global section} of \(E\) is a continuous map \(\sigma\colon M \to E\) such that \(\pi \circ \sigma = \mathrm{id}_M\). A local section is a continuous map defined on an open set \(U \subset M\), \(\sigma \colon U \to M\) such that \(\pi \circ \sigma =\mathrm{id}_U\). \(\mathcal{S}\)-sections are sections that are \(\mathcal{S}\)-morphisms. \index{Section of a vector bundle}}

Since we will usually work with \(\mathcal{S}\)-sections, the word section will mean \(\mathcal{S}\)-section unless otherwise stated. Following the notation in \parencite{lee_introSM}, a \textit{rough section}\index{Section!rough} is a section that needs not be continuous. We will denote \(\Gamma(M;E)\) and \(\Gamma(U;E)\) the set of global and local smooth sections on a smooth manifold, and \(\mathcal{O}(U;E)\) and \(\mathcal{O}(M;E)\) the corresponding sets for holomorphic sections. We make the distinction now because we will need it later. However, for the basic definitions and properties, there is no difference between smooth and holomorphic sections, so we will use \(\mathscr{E}(M;E)\) and \(\mathscr{E}(U;E)\) to denote either of them. We will use the notations \(\sigma(p)\) or \(\sigma_p\) interchangeably to denote the value of the section at \(p \in M\).

% \begin{example} \(\hspace{1 pt}\)
%     \begin{enumerate}
%         \item \(\mathfrak{X}(M)=\Gamma(M; TM)\) for any smooth manifold \(M\).
%         \item \(\mathcal{S}(M)=\mathscr{E}(M; M \times \K^r)\). Indeed, an \(\mathcal{S}\)-morfism \(f\colon M \to \K^r\) can be identified with a section on the trivial bundle \(M \times \K^r\). An \(\mathcal{S}\)-section on \(M \times \K^r\) is given by an \(\mathcal{S}\)-morphism \(x \mapsto (x, f(x))\), and it can be identified with the map \(f\).
%     \end{enumerate}
% \end{example}

Observe that \(\mathscr{E}(M;E)\) is a \(\K\)-vector space under pointwise addition and scalar multiplication: if \(\lambda, \mu \in \K\) and \(\sigma_1, \sigma_2 \in \mathscr{E}(M;E)\), then we can define
\[
(\lambda \sigma_1+\mu\sigma_2)(p):=\lambda \sigma_1(p) + \mu \sigma_2(p).
\] 
The sum takes place on each \(E_p\), so it is well-defined, and the result is still a section. Observe that we get an \(\mathcal{S}\)-section. Indeed, the map is an \(\mathcal{S}\)-morphism because we may always compose it with a trivialization \((U, \Phi)\). In that case, since \(\Phi\) is linear isomorphism on each \(E_p\), we get
\[
\Phi(\lambda \sigma_1+\mu\sigma_2)=(\mathrm{id}_U, \lambda \Phi(\sigma_1)+\mu \Phi(\sigma_2)).
\]
This is an \(\mathcal{S}\)-morphism because each component is, and since \(\Phi\) is an \(\mathcal{S}\)-isomorphism, we conclude that \(\lambda \sigma_1+\mu\sigma_2\) is an \(\mathcal{S}\)-morphism as well. Similarly, we can show that \(\mathscr{E}(M; E)\) is an \(\mathcal{S}(M)\)-module. Again, the product of a function \(f \in \mathcal{S}(M)\) and an \(\mathcal{S}\)-section is defined pointwise; \((f \sigma)(p):=f(p) \cdot \sigma(p)\), and an analogous procedure to the one we just did shows the claim. \\

As happens with vector fields, it is very convenient to have a set of sections that span the vector space \(E_p\) at each point \(p\) on some open set. This motivates the definition of local and global frames.

\Definition{\textbf{Frame}}{Let \(\pi\colon E \to M\) be a rank \(r\) \(\mathcal{S}\)-bundle. A set of \(r\) global (resp. local on \(U\)) \(\mathcal{S}\)-sections \(s_1, \ldots s_r \in \mathscr{E}(M; E)\) form a \textit{global frame}\index{Frame!global} (resp. \textit{local frame}\index{Frame!local} on \(U\)) if for each \(p \in  M\) (resp. \(p \in U\)), \(\corch{s_1(p), \ldots s_r(p)}\) is a basis of \(E_p\).}

\begin{example} \label{example_local_frame_local_triv}
    \(\hspace{1 pt}\)
    \begin{enumerate}
        \item Let \(M\) be an \(\mathcal{S}\)-bundle, and consider a trivial bundle \(E=M \times \K^r\). The standard basis \(\corch{e_1, \ldots e_n}\) on \(\K^r\) yields a global frame as follows: define \(\tilde{e}_i\colon M \to M \times\K^r, p \mapsto (p, e_i)\). This clearly defines \(\mathcal{S}\)-morphism, and since \(E_p=\corch{p} \times \K^r\), \(\corch{\tilde{e}_1(p), \ldots \tilde{e}_n(p)}\) is a basis for each \(p \in M\).
        \item Let \(\pi \colon E \to M\) be a rank \(r\) \(\mathcal{S}\)-vector bundle, and let \((U, \Phi)\) be a local trivialization. Define \({s}_i:=\Phi^{-1} \circ \tilde{e}_i\) for \(1 \leq i \leq r\). Then, \(\corch{{s}_1, \ldots {s}_r}\) is a local frame for \(E\) over \(U\). 
    \end{enumerate}
\end{example}
The local frames defined in the previous example are all possible local frames, as shown by the following proposition.

\Proposition{\textbf{Local frames associated with local trivializations}}{Let \(\pi\colon E \to M\) be an \(\mathcal{S}\)-vector bundle. Then, there is a one-to-one correspondence between local frames and local trivializations.\label{prop_local_frames_associated_trivializations}}
\begin{proof}
    Given a local trivialization \((U, \Phi)\), we define sections as in example \ref{example_local_frame_local_triv}. Reciprocally, given a local frame \((s_1, \ldots s_r)\), a trivialization sends \(\lambda_1\cdot s_1(p) + \ldots+ \lambda_r \cdot s_r(p)\) to \((p, (\lambda_1, \ldots \lambda_r))\). Details may be found in \parencite[Prop. 10.19]{lee_introSM}
    % The maps \(\tilde{e}_i\) are \(\mathcal{S}\)-morphisms, therefore, so are the maps \(s_i\). The fact that they are sections follows from the fact that the following diagram is commutative.
    % \begin{center}
    %     \begin{tikzcd}[row sep=normal, column sep=normal, text height=1ex, text depth=1ex]
    %         \pi^{-1}(U) \arrow[rr, "{\displaystyle \Phi}"] \arrow[rd, "{\displaystyle \pi}"] && U \times \K^r \arrow[ld, "{\displaystyle \pi_U}", ']  \\
    %         & U \arrow[bend left = 30, ul, "{\displaystyle s_i}"] \arrow[bend right = 30, ur, "{\displaystyle \tilde{e}_i}", ']& 
    %     \end{tikzcd}
    % \end{center}
    % Indeed, \((\pi \circ s_i)(p) = (\pi\circ \Phi^{-1})(p, e_i)=\pi_U(p, e_i)=p\) by definition of local trivialization. Finally, for each \(p \in U\), 
    % \[
    % \corch{s_1(p), \ldots s_r(p)}=\parent{\Phi|_{E_p}}^{-1}\big( \corch{e_1, \ldots e_r} \big),
    % \]
    % so \(\corch{s_1(p), \ldots s_r(p)}\) is the image of a basis by a linear isomorphism. \\

    % Reciprocally, let \((s_1, \ldots s_r)\) be a local frame over an open domain \(U \subset M\). Define \(\Psi\colon U \times \K^r \to \pi^{-1}(U), (p, (v^1, \ldots v^r))\mapsto v^j \cdot s_j(p)\). Observe that \(s_i=\Psi\circ \tilde{e}_i\) This map is clearly bijective because \(\corch{s_1(p), \ldots s_r(p)}\) is a basis of each \(E_p\) and, if we show that it is a local \(\mathcal{S}\)-isomorphism, it will be an \(\mathcal{S}\)-isomorphism, and its inverse will be the trivialization we seek. To do so, let \(q \in U\), and choose another local trivialization \((U, \Phi)\), with \(q \in V \subset U\). Observe that \(\Phi\circ s_i(p)=(p, (s_i^1(p),\ldots s_i^r(p)))\) for some \(\mathcal{S}\)-morphisms \(s_i^k\), because the left-hand side is an \(\mathcal{S}\)-morphism. Thus, on \(V \times \K^r\), we have
    % \[
    % \Phi \circ \Psi(p, (v^1, \ldots , v^r))=(p,(v^i\cdot s_i^1(p), \ldots ,v^i\cdot s_i^r(p))),
    % \]
    % which is an \(\mathcal{S}\) because so are its components. \((\Phi \circ \Psi)^{-1}\) is smooth because
    % \[
    % (\Psi \circ \Phi)^{-1}(p, (w^1, \ldots , w^r)) = (p,(w^i\cdot \tau_i^1(p), \ldots ,w^i\cdot \tau_i^r(p))),
    % \]
    % where \(\tau^j_k\) are the coefficients of the inverse of \(\sigma^\ell_m\), thus \(\mathcal{S}\)-morphisms.
\end{proof}

% The previous proposition allows us to obtain charts for vector bundles once we know a local frame. Those charts are in fact quite reasonable. Suppose that \((V, \varphi)\) is a chart on some neighborhood of a point \(p \in M\), and let \(\corch{\sigma_1, \ldots \sigma_r}\) be a local frame on \(V\). An element \(\xi \in E\) will lie on some \(E_p\), so we may write it as \(\xi=v^i\sigma_k(p)\). Then the "reasonable coordinates" for this point are \(\varphi(p), (v^1, \ldots v^r)\). This is precisely what the following corollary states.
% \begin{corollary}
%     Let \(\pi\colon E \to M\) be a rank \(r\) \(\mathcal{S}\)-vector bundle over \(M\), and \((V, \varphi)\) a smooth chart. Suppose that \(\corch{\sigma_1, \ldots \sigma_r}\) is a local frame on \(V\), and define \(\tilde{\varphi}\colon \pi^{-1}(V) \to \varphi(V) \times \K^r\) by
%     \[
%     \tilde{\varphi} \parent{v^i\sigma_i(p)} = (\varphi(p), (v^1, \ldots v^r)).
%     \]
%     Then, \(\tilde{\varphi}\) is a chart for \(E\).
% \end{corollary}
% \begin{proof}
%     It is enough to show that \(\tilde{\varphi}\) is an \(\mathcal{S}\)-isomorphism onto its image. Let \((V, \Phi)\) be local trivialization associated to the local frame. Then, \(\tilde{\varphi} = (\varphi \times \mathrm{id}_{\K^r}) \circ \Phi\) is a composition of \(\mathcal{S}\)-isomorphisms. 
% \end{proof}

% These results allow us to generalize the fact that a vector field is smooth if and only if its coordinate functions with respect to any chart (more precisely, with respect to the vector fields \(\corch{\partial_1, \ldots \partial_n}\)) are smooth. More precisely, if \(\corch{s_1, \ldots s_r}\) is a local frame of an \(\mathcal{S}\)-vector bundle over some open set, then for every local section \(\tau\colon V \to E\) there are functions \(\tau^1, \ldots \tau^r\) such that \(\tau=\tau^is_i\). This is obvious pointwise due to the fact that \(\corch{s_1(p), \ldots s_r(p)}\) form a basis of \(E_p\) for each \(p \in V\). So we are only concerned with studying whether those maps are \(\mathcal{S}\)-morphisms.

% \Proposition{\textbf{Local frame Criterion for \(\mathcal{S}\)-sections}}{Let \(\pi\colon E \to M\) be an \(\mathcal{S}\)-vector bundle, and \(\tau\colon M \to E\) a rough section. Let \(\corch{s_1, \ldots s_r}\) be an \(\mathcal{S}\)-frame for \(E\) over an open set \(U \subset M\). Then, \(\tau\) is an \(\mathcal{S}\)-section on \(U\) if and only if each \(\tau^i\colon U \to \K\) is an \(\mathcal{S}\)-morphism. }
% \begin{proof}
%     Let \((U, \Phi)\) be the local trivialization associated with the frame. Then, \(\Phi \circ \tau\colon U \to U \times \K^r\) satisfies \(\Phi \circ \tau = (\mathrm{id}_U, \tau^1, \ldots \tau^r)\). Since \(\Phi\) is an \(\mathcal{S}\)-isomorphism, \(\tau\) is an \(\mathcal{S}\)-morphism if and only if \((\mathrm{id}_u, \tau^1, \ldots \tau^r)\) is an \(\mathcal{S}\)-morphism, if and only of \(\tau^1, \ldots \tau^r\) are \(\mathcal{S}\)-morphisms.
% \end{proof}
Changes of coordiantes with respect to local trivializations induce changes of coordinates with respect to local frames as follows.
\Proposition{}{Let \((s_1, \ldots s_r)\) and \((\tilde{s}_1, \ldots \tilde{s}_r)\) be two \(\mathcal{S}\)-local frames defined on overlapping open domains. Then, there are \(\mathcal{S}\)-morphisms \(\tau^j_k\) such that \(\tilde{s}_k=\tau^j_k \cdot s_j\) and the matrix \(\tau=(\tau^j_k)_{1 \leq j,k\leq r}\) is non-singular. \label{prop_change_of_coordintes_frames}}
\begin{proof}
    By proposition \ref{prop_local_frames_associated_trivializations}, the frames \((s_1, \ldots s_r)\) and \((\tilde{s}_1, \ldots \tilde{s}_r)\) are associated to local trivializations \((U, \Psi)\) and \((V, \Phi)\) respectively. \(U\) and \(V\) are the definition domains of the frames, so \(U \cap V \neq \emptyset\). By lemma \parencite[Lemma 10.5]{lee_introSM} there is an \(\mathcal{S}\)-morphism \(\tau \colon U \cap V \to \mathrm{GL}(k, \K)\) such that \(\Phi\circ\Psi^{-1}(p,v)=(p, \tau(p)v)\). Since each frame defines a basis for each fiber, there are functions \(\lambda^j_k\) satisfying \(\hat{s}_k=\lambda^j_k \cdot s_j\). Now, applying the trivializations, we get
    \[
    \tilde{e}_k=(\mathrm{id}_{U \cap V}, e_k)=\Psi(\tilde{s}_k)=\Psi(\lambda^j_k\cdot s_j)= \Psi \circ \Phi^{-1}(\mathrm{id}_{U \cap V},\lambda^j_k \cdot e_j)=(\mathrm{id}_{U \cap V}, \lambda^j_k \cdot \tau^{-1}\cdot e_j)
    \]
    hence, \(\lambda^j_k \cdot \parent{\tau^{-1}}^i_j=\delta^i_k\), so \(\lambda^j_k=\tau^j_k\), which is an \(\mathcal{S}\)-morphism.
\end{proof}

Finally, let us study frames on the dual bundle \(E^*\) of an \(\mathcal{S}\)-bundle \(E\). Recall that every basis \(\corch{e_1, \ldots e_n}\) of a finite-dimensional vector space admits a dual basis \(\corch{e_1^*, \ldots e_n^*}\) defined by \(e_i^*(e_j)=\delta_{ij}\). We seek a similar construction for vector bundles: given a local frame \(\corch{s_1, \ldots s_r}\) on \(E\), we look for a dual frame \(\corch{t_1, \ldots t_n}\) on \(E^*\) such that \(t^j(s_i)=\delta^j_i\). Of course, the fact that the sections \(t^j\) exist is clear, we can just define rough sections that define the dual basis to \(\corch{s_1(p), \ldots s_r(p)}\) at each fiber. %It then remains to show that the \(t^j\)'s are \(\mathcal{S}\)-sections. {\color{red}terminar}.


\subsection{Bundle homomorphisms}
% Once defined the structure we are interested in, we look for the right definition of a morphism between those objects. There are two things that seem reasonable to ask. The first one is that a morphism should respect the \(\mathcal{S}\)-structure of the \(\mathcal{S}\)-bundles, so it should be an \(\mathcal{S}\)-morphism. Secondly, the core idea of vector bundles is to attach a vector point to each point on the manifold, so we would like to respect that as well. That is, vector bundle morphisms should restrict to linear maps between fibers. With this motivation in mind, we give the following definition.
Since the key idea of vector bundles is to attach a vector space to each point of a manifold, bundle morphisms to respect that linear structure, i.e. restrict to linear maps between fibers.

\Definition{\textbf{Vector bundle morphism}}{Let \(\pi\colon E \to M\) and \(\pi'\colon E' \to M\) be two \(\mathcal{S}\)-vector bundles over an \(\mathcal{S}\)-manifold \(M\). A morphism of \(\mathcal{S}\)-vector bundles, \(F\), is an \(\mathcal{S}\)-morphism \(F\colon E \to E'\) satisfying: \index{Vector bundle!isomorphism} \label{def_vector_bundle_morphism}

\begin{minipage}{.5\textwidth}
    \begin{enumerate}
        \item \(\pi' \circ F = \pi\)
        \item \(F|_{E_p}\colon E_p \to E'_{F(p)}\) is a linear map.
    \end{enumerate}
    \end{minipage}\hfill 
    \begin{minipage}{0.5\textwidth}
        \centering
        \begin{tikzcd}[row sep=normal, column sep=normal,  ampersand replacement=\&]
            E \arrow[rr, "{\displaystyle F}"] \arrow[rd, "{\displaystyle \pi}", ']  \& \& E' \arrow[ld, "{\displaystyle \pi'}"]  \\
            \& M \& 
        \end{tikzcd}
    \end{minipage}
}
Of course, we will say that \(F\colon E \to E'\) is an \(\mathcal{S}\)-vector bundle isomorphism if it is bijective and its inverse also is a bundle morphism. In that case, we say that \(E\) and \(E'\) are isomorphic. 

% \Proposition{}{Let \(\pi\colon E \to M\) and \(\pi'\colon E' \to M\) be two \(\mathcal{S}\)-vector bundle isomorphisms, and 
% \(F\colon E \to E'\) a vector bundle morphism. If \(F\) is bijective, then it is an isomorphism. \label{bundle_isomorphims_bijection}}
% \begin{proof}
%     {\color{red}Añadir si lo utilizo}
% \end{proof}

\begin{obs}
    \(\hspace{1 pt}\)
    \begin{enumerate}
        \item Given a vector bundle \(E\), the identity map \(\mathrm{id}_E\colon E \to E\) is trivially a vector bundle morphism (an isomorphism in fact).
        \item Observe that an \(\mathcal{S}\)-bundle over \(M\) is trivial if and only if it is isomorphic to \(M \times \K^r\). Indeed, an isomorphism \(F\colon M \times \K^r \to E\) is a global trivialization. Just observe that \(F\) is bijective because it is an isomorphism, and it is an \(\mathcal{S}\)-morphism and commutes with the projections by definition \ref{def_vector_bundle_morphism}, so it satisfies all the properties of the definition of a global trivialization. Reciprocally, any global trivialization will define a vector bundle isomorphism using the same reasoning.
        \item The set of bundle homomorphisms \(E \to E'\) can be identified with the space of sections of the vector bundle \(\mathrm{Hom}(E,E')\). Indeed, each bundle morphism \(F\colon E \to E'\) defines a section \(\hat{F}\in \mathscr{E}(M;\mathrm{Hom}(E,E'))\) by setting \(\hat{F}_p:=F|_{E_p}\). Reciprocally, a section \(\hat{F}\) defines a vector bundle morphism \(F\) by setting \(F(v):=\hat{F}_p(v)\) if \(v \in E_p\). As a particular case of this construction, there is a vector bundle isomorphism \(\mathrm{End}(E)\cong E^* \otimes E\).
    \end{enumerate}
\end{obs}

\begin{example}
    \(\hspace{1 cm}\)
    \begin{enumerate}
        \item A well-known example comes from riemannian geometry: if \(M\) is a riemannian geometry, the \textit{musical isomorphims} define a vector bundle isomorphism between \(TM\) and \(T^*M\). Analogously, a symplectic form defines a bundle isomorphim between the same vector bundles.
        \item Given two \(\mathcal{S}\)-vector bundles \(E,E'\) over the same \(\mathcal{S}\)-manifold, \(M\), we have an \(\mathcal{S}\)-bundle isomorphism \(\mathrm{Hom}(E,E')\cong E^*\otimes E'\). This eventually comes down to the fact that, for every pair of finite-dimensional \(\K\)-vector spaces \(V\) and \(W\), there is an isomorphism \(\mathrm{Hom}(V,W) \cong V^* \otimes W\). Indeed, define \(\Phi\colon V^*\times W \to \mathrm{Hom}_\K(V, W), (\alpha,w)\mapsto \phi_\alpha^w\), where \(\phi_\alpha ^w(v):=\alpha(v)\cdot w \in W\). \(\Phi\) is clearly bilinear, so by the universal property of the tensor product, it descends to a linear map \(\hat{\Phi}\colon V^* \otimes W \to \mathrm{Hom}_\K(V, W)\). Since both spaces are finite-dimensional, \(\dim_\K (V^* \otimes W)=\dim_\K (\operatorname{Hom}_\K(V,W))\), so it suffices to show that \(\hat{\Phi}\) is surjective. Indeed, choose a basis \(\corch{e_1, \ldots e_n}\) of \(V\), and take its dual basis \(\corch{e_1^*, \ldots e_n^*}\). Choose \(\phi \in \operatorname{Hom}_\K(V, W)\), and let \(w_i:=\phi(e_i)\). Then, \(\phi=\hat{\Phi}(\sum_{i=1}^n e_i^* \otimes w_i)\).
    \end{enumerate}
\end{example}

% \begin{example}
%     There is an example of bundle isomorphisms that arises naturally in the context of symplectic geometry which we will use extenseively. Let \((M, \omega)\) be a symplectic manifold. We are going to define an analogue to the \textit{musical isomorphisms}\index{musical isomorphisms} of riemannian geometry in the symplectic setting (in fact, the isomorphisms we are going to define are also called musical isomorphisms). We want to define a bundle isomorphism \(\omega^\flat\colon TM \to T^*M\), so we will define its action on the fibers first. We thus define
%     \[
%     \begin{array}{cccc}
%         \omega^\flat_p\colon & T_pM & \longrightarrow & T_p^*M \\
%         & v & \mapsto &\omega^\flat_p(v),
%     \end{array}
%     \]
%     where \(\omega_p^\flat(v)(w):=\omega_p(v,w)\). Since \(\omega_p \in \Lambda^2(T_pM)\), it is multilinear, so \(\omega_p^\flat\) is linear. Moreover, we may glue all those maps to define \(\omega^\flat\colon TM \to T^*M\) by \(\omega^\flat(v)(w):=\omega_p(v,w)\) for \(v,w \in T_pM\). Observe that for each \(p \in M\), \(\omega^\flat_p\) is bijective. Indeed, since \(\dim T_pM=\dim T_p^*M\) we just need to show that it is injective, but this is obvious because \(\omega\) is a symplectic form (see remark \ref{dim_sympl_v_spaces}). This means that \(\omega^\flat\) is bijective, so by proposition \ref{bundle_isomorphims_bijection}, we only need to check that \(\omega^\flat\) is smooth (note that \(\omega^\flat\) trivially commutes with the standard projections). By the vector bundle characterization lemma {\color{red} añadir}, we only need to check that whenever \(X \in \mathfrak{X}(M)\), \(\omega^\flat(X)\) is \(\mathscr{C}^\infty(M; \R)\), which is obvious because \(\omega^\flat\) acts linearly fiberwise. Therefore, \(\omega^\flat\) is an isomorphims between \(TM\) and \(T^*M\), and has an inverse \(\omega^\sharp\colon T^*M \to TM\). \\

%     The previous isomorphisms are usually expressed in terms of \(\mathscr{C}^\infty(M; \R)\)-linear maps between the spaces of sections: \(\mathfrak{X}(M)\) and \(\Omega^1(M)\) That is, they can be used to assign a \(1\)-form to each vector field and vice versa. In particular, \(\omega^\flat(X)=i_X\omega \in \Omega^1(M;\R)\). 
% \end{example}

\subsection{Subbundles}
Subbundles are a way to formalize the natural notion of "taking a subspace of a vector bundle" that one would come up with. More precisely, if we are to take a reasonable subset of a vector bundle, and we want to still call it a bundle, the first reasonable think to ask is that the intersection of the subset we have taken and each of the fibers of the original vector bundle is a vector subspace of the fiber. Since we want to make a bundle out of it, we must ask for all the so defined vector subspaces so have the same dimension. Furthermore, since we work with vector bundles that are \(\mathcal{S}\)-manifolds, we should ask subbundles to be embedded \(\mathcal{S}\)-submanifolds.

\Definition{\textbf{Subbundle}}{Let \(\pi\colon E \to M\) be an \(\mathcal{S}\)-bundle over an \(\mathcal{S}\)-manifold \(M\). An \(\mathcal{S}\)-subbundle is an embedded submanifold \(D \subset E\) such that \(D_p:=D \cap E_p\) is a vector subspace of \(E_p\) for each \(p \in M\), all of them of the same dimension, together with the restriction \(\pi|_D\colon D \to M\), such that \(\pi|_D\colon D \to M\) is a vector bundle.}

The following provides a very useful way to recognize subbundles.

\Lemma{\textbf{Local frame criterion for subbundles}}{Let \(\pi\colon E\to M\) be a rank \(r\), \(\mathcal{S}\)-bundle, and suppose that, for each \(m \in M\), we are given an \(\ell\)-dimension subspace \(D_m \subset E_m\). Then, \(D=\cup_{m \in M} D_m \subset E\) is an \(\mathcal{S}\)-subbundle if and only if each point \(m \in M\) has a neighborhood \(U\) on which there are \(\mathcal{S}\)-local sections \(\corch{\sigma_1, \ldots \sigma_\ell}\) such that \(\corch{\sigma_1(m), \ldots \sigma_\ell(p)}\) is a basis of \(D_m\). \label{lemma_local_frame_crit}
}
\begin{proof}
    See \parencite[Lemma 10.32]{lee_introSM} for a proof in the smooth category, which can be adpated verbatim to the holomorphic category.
\end{proof}

% \subsection{Some useful identifications}
% When working with vector bundles and their sections, there are several identifications that allow us to view some objects as other related objects. It is important to keep track of the identifications we are going to present, because in the literate, they are often used without explicit warning. The following theorem gives a summary of some of the identifications we will care about.

% \Theorem{}{Let \(E\) and \(F\) be two smooth vector bundles. Then there are bijections between the following sets:
% \begin{enumerate}
%     \item The set of \(\mathscr{C}^\infty(M;\R)\)-linear maps \(\Gamma(E)\to \Gamma(F)\)
%     \item The set of smooth bundle homomorphisms \(E \to F\)
%     \item The set of smooth sections of \(\mathrm{Hom}(E,F)\)
%     \item The set of smooth sections of \(E^*\otimes F\).
% \end{enumerate}
% }
% \begin{proof}
%     We will not give a detailed proof of this result, but we will write down the bijections explicitly. {\color{red} Terminar.}
% \end{proof}


\section{The complexified Tangent and Cotangent Bundles} \label{section_complexified_bundles}
Observe that under the identification \(\R^{2n} \equiv \C^n\), every complex manifold is a smooth manifold. Hence, it makes sense to speak about vector fields and covector fields (one-forms) in any complex manifold. However, we would like to extend those definitions to complex vector fields and one-forms (we will introduce higher-order complex differential forms later on).

\subsection{Complexification of real vector spaces} \label{section_complexification}
Firstly, given a real vector space \(V\), we may always consider its \textit{complexification}\index{Complexification!of a vector space}, which is the complex vector space \(V_\C:=V \otimes_\R \C\). Its elements can always be written as \(u\otimes 1+ v \otimes \ii\), and we shall omit the tensor product and simply write \(u+ \ii v \in V_\C\). Every \(\R\)-linear map \(f\colon V \to W\) defines a \(\C\)-lineal map \(f_\C\colon V_\C \mapsto W_\C\) (which we will usualy denote \(f\) as well) by linear extension of \(f_\C(u \otimes \alpha):=(f(u)) \otimes \alpha\) for \(\alpha \in \C\). Complexifications are endowed with natural conjugation maps \(\overline{u+ \ii v}:=u-\ii v\). This allows us to speak about real vectors, which are those \(w \in V_\C\) for which \(\overline{w}=w\). Obsserve that if we set \((V_\C)^\R\) the underlying real subspace, then \((V_\C)^\R\neq V\) (for example, \(\R_\C\cong \C\), but \(\C^\R\cong \R^2\)). Moreover, one has \((V_\C)^\R=V\oplus \ii V\). We usually write \(V_\C=V \oplus \ii V\) making some abuse of notation, but the right-hand side is not a direct sum decomposition of the left-hand side as complex vector spaces.

\subsection{Complexification of real vector bundles}
Let \(E\) be a real vector bundle over a smooth manifold \(M\). The \textit{complexification}\index{Complexification!of a vector bundle} of \(E\) is the vector bundle \(E_\C:=E \otimes  \C\). More precisely, we consider the tensor product of the real vector bundle \(E\) and the trivial bundle \(M \times \C\) (seen as a real vector bundle). Then, \(E \otimes (M \times \C)\) defined as above is a real vector bundle whose fiber at each \(p \in M\) is just \(E_p \otimes_\R \C\), that is, \((E_p)_\C\). Thus, \(E_\C\) is endowed with the structure of a complex vector bundle. Observe that it is smooth, but not necessarily holomorphic. To sum up, \(E_\C\) is the complex vector bundle whose fiber at each point is the complexification of the fiber of the original bundle (see \parencite[Ch. 1]{lee_introCM} for a more explicit construction). \\

The \textit{complexified tangent and cotangent bundles} are the complexifications of \(TM\) and \(T^*M\) respectively. They are denoted \(T_\C M\) and \(T_\C^* M\) respectively. Complex vector and covector fields are just sections of these bundles. Observe that \(T_\C M\) can be interpreted as a very natural object. Recall from basic differential geometry, that for a smooth manifold, \(T_pM\) is the vector space of derivations acting on (germs of) smooth real-valued maps. Then, there is a complex vector space isomorphism between \(T_pM \otimes_\R \C\) and the space of \(\C\)-linear derivations acting on germs of complex-valued smooth maps {\color{red} revisar}. What that means for practical purposes is that, given \(Z=X+\ii Y \in \Gamma(M; T_\C M)\) and \(f=u+\ii v \in \mathscr{C}^\infty(M;\C)\), we have 
\[
Z(f)=\parent{X(u)-Y(v)} + \ii \parent{Y(u)+X(v)}.
\]
We can define a Lie bracket on \(\Gamma(M; T_\C M)\) by complex-linear extension of the original Lie bracket on \(TM\):
\[
[X_1+\ii Y_1, X_2+ \ii Y_2]:=\big([X_1, X_2]-[Y_1, Y_2] \big) + \ii \big([Y_1, X_2]+[X_1, Y_2]\big),
\]
which turns \((\Gamma(M; T_\C M), [\cdot, \cdot])\) into a Lie algebra. \\

Similarly, elements of \(\Gamma(M; T_\C^*M)\) are called complex one-forms. Since \(V_\C^*\) is naturally isomorphic to \((V^*)_\C\) (see \parencite[Lemma 1.2.6]{huybrechts2005complex}), for each \(p \in M\), we have that the fiber of \(T_\C^*M\) at \(p\), \(T_p^*M \otimes _\R \C\) is naturally isomorphic to \((T_pM \otimes_\R  \C)^*\). In the language of vector bundles, this means that \(T_\C^* M \cong \mathrm{Hom}_\C(T_\C M, \C)\). Then the action of a complex one form \(\alpha + \ii \beta\) on a complex vector field \(X + \ii Y\) is just given by \((\alpha + \ii \beta)(X+\ii Y)=(\alpha(X)-\beta(Y))+\ii (\alpha(Y)+\beta(X))\). Any \(f \in \mathscr{C}^\infty(M;\C)\) defines now a complex one-form \(\diff f:=\diff u + \ii \diff v\) (we will generalize the exterior differential when we introduce higher order complex differential forms). 

\section{Complex coordinate frames}
Let \(M\) be a complex manifold and choose holomorphic coordinates \((z^1, \ldots z^n)\). Each coordinate map \(z^j\) can be decomposed into its real and imaginary part \(x^j\) and \(y^j\), thus we get real coordinates \((x^1, y^1, \ldots, x^n, y^n)\). This yields a local smooth frame of \(TM\) given by \(\corch{\partial/\partial x^1, \partial/\partial y^1, \ldots \partial/\partial x^n, \partial/\partial y^n}\). Defining \(\partial/\partial z^j\) and \(\partial/\partial\bar{z}^j\) as in equation \eqref{derparz}, we obtain a new local frame for \(T_\C M\) which is given by
\[
\bigg\lbrace{\derpar{}{z^1}, \derpar{}{\bar{z}^1}, \ldots \derpar{}{z^n}, \derpar{}{\bar{z}^n} \bigg\rbrace}.
\]
A function \(f \in \mathscr{C}^\infty(M;\C)\) is holomorphic if and only if \(\partial f /\partial \bar{z}^j=0\) and antiholomorphic if and only if \(\partial f/\partial z^j=0\) for \(1 \leq j \leq n\) (see \parencite[Prop. 1.4.2]{lee_introCM}). Analogously, a local coframe for the complexified cotangent bundle is given by \(\corch{\diff z^1, \ldots \diff z^n, \diff \bar{z}^1, \ldots \diff \bar{z}^n}\), where we have defined (\(\diff x^j\) and \(\diff y^j\) are extended by complex linearity)
\[
\diff z^j:= \diff x^j+\ii \diff y^j \hspace{1 cm}
\diff \bar{z}^j := \diff x^j- \ii \diff y^j \hspace{1 cm} 1 \leq j \leq n.
\]
Then, we may express \(\diff f \in \Gamma(M; T_\C^*M)\) in either frame
\begin{equation} \label{exp_diff}
\diff f=\derpar{f}{x^j} \cdot \diff x^j + \derpar{f}{y^j} \cdot \diff z^j = \derpar{f}{z^j}\cdot \diff z^j + \derpar{f}{\bar{z}^j}\cdot \diff \bar{z}^j.
\end{equation}
Observe that the complexified bundles may be defined for smooth manifolds as well, but the local frames defined above only make sense for complex manifolds (because they require holomorphic coordinates). \\

Finally, consider two complex manifolds, \(M\) and \(N\), and let \(F \colon M \to N\) be a smooth map. From real differential geometry, we know that we mat consider its differential at a point \(m \in M\), which is an \(\R\)-linear map \(\diff_mF\colon T_mM \to T_{F(m)}N\). This map can be extended by complexification \(\diff_m F_\C\colon (T_mM)_\C \to (T_{F(m)}N)_\C\). We will usually denote it by the same symbol \(\diff_m F\). If we now write \(D'F(p)=(\partial F^k/\partial z^j)\), holomorphicity is characterized by the fact that, holomorphic map are those whose differential has the following matrix form
\[
\begin{pmatrix}
    D'F(p) & 0 \\
    0 & \overline{D'F(p)}
\end{pmatrix}
\]
with respect to any pair of holomorphic frames (associated to holomorphic coordinates of the manifold). See \parencite[Ch. 1]{lee_introCM} for more details.

\section{Almost complex structures} \label{appendix_almost_complex_structure}
Let \(V\) be a real vector space. An endomorphism \(J\colon V \to V\) is called an \textit{almost complex structure}\index{Almost complex structure} if \(J^2:= J \circ J=-\mathrm{id}_V\). Almost complex structures receive their name because, when they exist, they endow \(V\) with the structure of a complex vector space. Indeed, for \(a+\ii b \in \C\) and \(v \in V\) their product can be defined as \((a+ \ii b)*v:=a \cdot  v+ b\cdot  J(v)\). As seen in proposition \ref{almost_compl_decom} bellow, only even-dimensional real vector spaces can have almost complex structures. \\

Given a real vector space endowed with an almost complex structure, the fact that \(J^2=-\mathrm{id}_V\) means that \(J\) has no real eigenvalues. If we extend \(J\) to the complexified space, then the only possible eigenvalues of \(J\) are \(\ii\) and \(-\ii\). If we use the symbol \(J\) acting on vectors of a complexified vector space, we will always understand that we are working with the complexified map.

\Proposition{}{Let \(J\) be a complex structure on a real vector space \(V\). Then \(V_\C\) has a complete eigenspace decomposition \(V_\C=V^{(1,0)}\oplus V^{(0,1)}\), where \(V^{(1,0)}\) is the \(\ii\)-eigenspace and \(V^{(0,1)}\) is the \((-\ii)\)-eigenspace. Moreover, the decomposition is given, for \(w \in V_\C\) by 
\[
w^{(1,0)}=\frac{1}{2}\cdot \parent{w-\ii J(w)} \hspace{1 cm}
w^{(0,1)}=\frac{1}{2}\cdot (w+\ii J(w)).
\]
If \(\dim_\R V<\infty\), then \(\dim_\C V^{(1,0)}=\dim_\C V^{(0,1)}\), and therefore, \(\dim_\C V_\C=\dim_\R V\) is even.  \label{almost_compl_decom}   
}
\begin{proof}
    See \parencite[Prop. 1.52]{lee_introCM}. 
\end{proof}

\begin{obs}
    Let \(V\) be an even-dimensional real vector space endowed with an almost complex structure \(J_V\). We have seen that this induces on \(V\) the structure of a complex vector space. On the other hand, the complexification \(V_\C\) induces another linear complex structure on \(V\). \textit{These structures are not the same!} Indeed, \(\dim_\C V_\C=\dim_\R V=2\cdot \dim_\C(V,J_V)\). Another way of seeing this is that for \(v\in V^{(0,1)}\) \(J_V(v)\) and \(\ii v\) are opposites, even though both of them can be thought of as ``multiplying by the imaginary unit''.
\end{obs}

These constructions can be applied to vector bundles. Given a real vector bundle \(E\) over a smooth manifold, a vector bundle endomorphism \(J \colon E \to E\) is an \textit{almost complex structure} if \(J^2=-\mathrm{id}_E\).

\begin{obs}
    Observe that \(\C^n\) admits an almost complex structure. Indeed, \(T\C^n\) (seen as a real vector bundle) has an almost complex structure, which is given by the endomorphism
    \begin{equation} \label{almost_com_struc}
    J\parent{\derpar{}{x^j}}=\derpar{}{y^j} \hspace{1 cm} J\parent{\derpar{}{y^j}}=-\derpar{}{x^j}.
    \end{equation}
\end{obs}
Every complex manifold admits a complex structure on \(TM\), which is given, in terms of any holomorphic coordinate system \((z^1=x^1+\ii y^1, \ldots z^n=x^n+\ii y^n)\) by equation \eqref{almost_com_struc} (see \parencite[Prop. 1.55]{lee_introCM}). In that case, \(J\) induces a decomposition in the fibers of \(T_\C M\) by proposition \ref{almost_compl_decom}. This can be elevated to a vector bundle decomposition: indeed, for every complex manifold there is a Whitney sum decomposition (see example \ref{ex_vector_bundle_constructions})
\begin{equation} \label{decomp_T_CM}
T_\C M= T^{(1,0)}M \oplus T^{(0,1)}M,
\end{equation}
where the fibers at each point of \( T^{(1,0)}M\) and \(T^{(0,1)}M\) are the \(\ii\) and \((-\ii)\)-eigenspaces of \(J\) at each point. Moreover, local frames for those complex smooth subbundles are
\[
\bigg\lbrace \derpar{}{z^1}, \ldots \derpar{}{z^n} \bigg\rbrace \hspace{1 cm} \text{and} \hspace{1 cm}
\bigg\lbrace \derpar{}{\bar z^1}, \ldots \derpar{}{\bar z^n} \bigg\rbrace.
\]
Observe that, indeed we have
\[
J\parent{\derpar{}{z^j}}=J\parent{\frac{1}{2}\cdot \parent{\derpar{}{x^j}-\ii \derpar{}{y^j}}}=\frac{1}{2}\cdot \parent{\derpar{}{y^j}+\ii\derpar{}{x^j}}=\ii \derpar{}{z^j},
\]
and analogously, \(J(\partial/\partial \bar{z}^j)=-\ii \partial/\partial \bar{z}^j\). In that case, \(T^{(1,0)}M\) has the structure of a holomorphic vector bundle, but \(T^{(0,1)}M\) is only a smooth (complex) vector bundle. This justifies their names: \(T^{(1,0)}M\) is referred to as the \textit{holomorphic tangent bundle}\index{Tangent bundle!holomorphic} and \(T^{(0,1)}M\) is the \textit{antiholomorphic tangent bundle}\index{Tangent bundle!antiholomorphic}. Proofs of these claims may be found in \parencite[Prop. 156 and Example 3.37]{lee_introCM}. Finally, we may also consider \(T_JM\), which is the usual tangent bundle, but seen as a complex vector bundle through the complex structure \(J\). It is isomorphic as a holomorphic vector bundle to \(T^{(0,1)}M\).

\begin{obs} \label{holomorp_func_zero}
    Let \((z^1,\ldots z^n)\) be holomorphic coordinates. Every \(Z \in T^{(0,1)}M\) can be written as \(Z=Z^j \cdot \partial/\partial \bar{z}^j\). If \(f \in \mathcal{O}_M\), then \(Z(f)=Z^j(f)\cdot \partial f/\partial \bar{z}^j = 0\).
\end{obs}

The canonical almost complex structures on complex manifolds characterize holomorphicity of smooth maps as follows.

\Theorem{}{Let \(M, N\) be two complex manifolds, and let \(J_M\) and \(J_N\) be their canonical almost complex structures. Let \(F\colon M \to N\) be a smooth map. Then, \(F\) is holomorphic if and only if \(\diff F \circ J_M=J_N \circ \diff F\). \label{th_caracterization_holomoprh_J}}
\begin{proof}
    See \parencite[Prop. 1.55]{lee_introCM}.
\end{proof}

Given an even-dimensional smooth manifold \(M\) it makes sense to ask whether it admits a complex structure. It does if it is complex, but that is not a necessary condition: there are smooth manifolds that admit almost complex structures, but do not admit a complex manifold structure. Let \(M\) be a smooth manifold. We say that \(M\) is an \textit{almost complex manifold}\index{Almost complex manifold} if \(TM\) is endowed with an almost complex structure. The decomposition given in equation \eqref{decomp_T_CM} remains valid when \(M\) is a \(2n\)-smooth manifold (see \parencite[Lemma 1.60]{lee_introCM}). We say that an almost complex structure \(J\) on a smooth manifold is integrable if, for every \(X,Y \in \Gamma(M; T^{(1,0)}M)\), \([X,Y] \in \Gamma(M; T^{(1,0)}M)\). Then, there is a theorem by Newlander and Niremberg  that establishes that an almost complex structure comes from a complex manifold structure on \(M\) if and only if it is integrable. Equivalently, an almost complex structure \(J\) is integrable if and only if its \textit{Nijenhuis tensor} \(N_J\) vanishes, i.e. if the following tensor is identically zero:
\[
\begin{array}{cccc}
    N_J\colon& \Gamma(T_\C M) \otimes \Gamma(T_\C M) & \longrightarrow & \Gamma(T_\C M) \\
     & (X,Y) & \mapsto &  [X,Y]-[JX, JY] + J [JX, Y] + J[X, JY].
\end{array}
\]

\section{The Dolbeault complex} \label{sec_Dolbeault_complex}
An almost complex structure \(J\) on a vector space \(V\) induces an almost complex structure on its dual by \(\big[J(f)\big](v):=f(J(v))\), which induces a natural decomposition on \(V^*_\C=(V^*)^{(1,0)} \oplus (V^*)^{(0,1)}\) and there are isomorphisms \((V^*)^{(1,0)} \cong \parent{V^{(1,0)}}^*\) and \((V^*)^{(0,1)} \cong \parent{V^{(0,1)}}^*\). This decomposition carries over to the complexified cotangent bundle, and yields \(T_\C^*M =(T^*M)^{(1,0)} \oplus (T^*M)^{(0,1)} \cong (T^{(1,0)}M)^* \oplus (T^{(0,1)}M)^*\). Moreover, local frames for \((T^*M)^{(1,0)}\) and \((T^*M)^{(0,1)}\) are given by 
\begin{equation} \label{frames_of_compl_cotang}
\corch{\diff z^1, \ldots \diff z^n} \hspace{1 cm} \text{and} \hspace{1 cm} \corch{\diff \bar{z}^1, \ldots \diff \bar{z}^n}
\end{equation}
respectively. \\

Let \(M\) be a complex manifold we define 
\[
{\bigwedge}^{(p,q)}_\C M:=\parent{{\bigwedge}^p (T^{(1,0)}M)^*} \otimes_\C \parent{{\bigwedge}^q (T^{(0,1)} M)^*}.
\]
A complex \((p,q)\)-form is a section of this complex vector bundle. We will denote \(\Omega_\C^{(p,q)}(M):=\Gamma(M; \bigwedge^{(p,q)}_\C M)\). Its elements are called \textit{complex differential forms}. {\color{red} Revisar} Moreover, each \(\bigwedge^{(p,q)}_\C M\) can be seen in a canonical way as a vector subbundle of \(\bigwedge^{p+q} T^*_\C M \). This follows from the fact that, at the vector space level, (all definitions for real vector spaces endowed with an almost complex structure are the same) we have a canonical linear injection
\[
{\bigwedge}^p V^{(1,0)} \otimes_\C {\bigwedge}^q V^{(0,1)} \longrightarrow {\bigwedge}^{p+q} V_\C
\]
given by the linear extension of \(\alpha \otimes \beta \mapsto \alpha \wedge \beta\) (to do this more rigorously we would have to invoke the universal property of the tensor product). What is more, we actually have a vector bundle Whitney sum decomposition
\begin{equation} \label{direct_sum_pq_forms}
{\bigwedge}^k_\C M:={\bigwedge}^{k \ } T^*_\C M = \bigoplus_{p+q=k} \parent{{\bigwedge}^{(p,q)}_\C M}.
\end{equation}
By setting \(\Omega_\C^{(p,q)}(M):=\Gamma(M; {\bigwedge}^{(p,q)}_\C M)\), whose elements are called \((p,q)\)\textit{-forms}\index{\((p,q)\)-forms} or forms of \textit{bidegree}\index{Bidegree of a complex form} \((p,q)\), the Whitney sum decomposition carries over to a Whitney sum decomposition of the spaces of complex differential forms 
\[
\Omega_\C^k(M)=\bigoplus_{p+q=k} \Omega_\C^{(p,q)}(M).
\]
There is another canonical isomorphism \({\bigwedge}^k_\C M \cong ({\bigwedge}^k T^*M) \otimes_\R \C\), which means that every complex \(k\)-form can be seen as a sum \(\alpha + \ii \beta\), where \(\alpha, \beta \in \Omega_\R^k(M)\). In particular, this means that there is a conjugation map in the space of complex differential forms, which is given by \(\overline{\alpha + \ii \beta}:=\alpha - \ii \beta\). By convention, we set \(\Omega_\C^{(p,q)}(M)=\corch{0}\) except if \(0 \leq p,q \leq n\). We define the graded algebra of differential forms by \(\Omega_\C^*(M):=\oplus _{k\geq 0} \Omega_\C^k(M)\). Observe that we also have \(\Omega_\C^k(M)=\Omega_\R^k(M) \otimes _\R \C\). \\

Observe that every vector space can be embedded into its complexification as a real vector subspace \(V \hookrightarrow V_\C, v \mapsto v \otimes 1\). In particular, we may see \(\bigwedge^k T^*M\) as a real subbundle of \(\bigwedge^k T_\C^* M =\bigwedge^k_\C M\). This means that we can see \(\Omega_\R^k(M)\) as a real vector subspace of \(\Omega_\C^k(M)\) by simply extending real forms to complex vectors by complex bilinearity.
\begin{obs} \label{charac_pq_forms}
    We could relax our hypothesis and ask \(M\) to be almost complex, but since we will work almost exclusively with complex manifolds, the added generality is pointless. Moreover, when \(M\) is a complex manifold, we have holomorphic coordinates \((z^1, \ldots z^n)\) the local coframes of equation \eqref{frames_of_compl_cotang} are available. More importantly, \((p,q)\)-forms are characterized by the fact that they can be written as
    \begin{equation} \label{coordinates_complex_diff_form}
    \omega = {\sum_{I \bar{J}}}' \omega_{I, \bar{J}} \cdot \diff z^{i_1} \wedge \ldots \wedge \diff z^{i_p} \wedge \diff \bar{z}^{j_1} \wedge \ldots \wedge \diff \bar{z}^{j_q}, 
    \end{equation}
    where \(I=(i_1, \ldots i_p)\) and \(\bar{J}=(j_1, \ldots j_q)\) are increasing sets of indices, and \(\omega_{I \bar{J}}\) are smooth complex-valued maps (see \parencite[Lemma 4.1]{lee_introCM}). The symbol \(\sum'\) means that we sum over ordered increasing sets of indices. Another useful characterization is that a complex \(p+q\)-form \(\omega\) is of type \((p,q)\) if and only if it vanishes whenever it is applied to more that \(p\) sections of \(T^{(1,0)}M\) or more than \(q\) sections of \(T^{(0,1)}M\).
\end{obs}
Decomposition \eqref{direct_sum_pq_forms} induces canonical projections
\[
\Pi^{(p,q)} \colon {\bigwedge}^{p+q}_\C M \longrightarrow {\bigwedge}^{(p,q)}_\C M.
\]
Moreover, the exterior differential \(\diff\) can be extended by complex linearity to \(\Omega_\C^k(M)\). Then, a simple coordinate calculation shows that \(\diff \parent{\Omega_\C^{(p,q)}(M)} \subset \Omega_\C^{(p+1, q)}(M) \oplus \Omega_\C^{(p, q+1)}(M)\) (see \parencite[Prop. 4.5]{lee_introCM})

% \Proposition{}{Let \(M\) be a complex manifold. Then \(\diff \parent{\Omega_\C^{(p,q)}(M)} \subset \Omega_\C^{(p+1, q)}(M) \oplus \Omega_\C^{(p, q+1)}(M)\).}
% \begin{proof}
%     Since \(\diff\) acts locally, let \((z^1, \ldots z^n)\) be local holomorphic coordinates and choose \(\alpha \in \Omega_\C^{(p,q)}(M)\). Since \(\diff^2=0\), we have
%     \begin{align*}
%         \diff \alpha & = \diff \parent{{\sum_{I \bar{J}}}' \alpha_{I, \bar{J}} \cdot \diff z^{i_1} \wedge \ldots \wedge \diff z^{i_p} \wedge \diff \bar{z}^{j_1} \wedge \ldots \wedge \diff \bar{z}^{j_q}} = {\sum_{I \bar{J}}}' \diff \alpha_{I, \bar{J}} \wedge \diff z^{i_1} \wedge \ldots \wedge \diff z^{i_p} \wedge \diff \bar{z}^{j_1} \wedge \ldots \wedge \diff \bar{z}^{j_q} \\
%         & = {\sum_{I \bar{J}}}' \parent{\derpar{\alpha_{I \bar{J}}}{z^k} \cdot \diff z^k + \derpar{\alpha_{I \bar{J}}}{\bar{z}^k} \cdot \diff \bar{z}^k} \wedge \diff z^{i_1} \wedge \ldots \wedge \diff z^{i_p} \wedge \diff \bar{z}^{j_1} \wedge \ldots \wedge \diff \bar{z}^{j_q}
%     \end{align*}
%     where we used \eqref{exp_diff} in the last equal sign. This expands to the sum of \((p+1, q)\) and \((p, q+1)\) forms so by the characterization above (see \eqref{coordinates_complex_diff_form}), we have the desired result.
% \end{proof}

It now makes sense to define the \textit{Dolbeault operators}\index{Dolbeault operators}
\[
\partial:=\Pi^{(p+1, q)} \circ \diff \colon \Omega_\C^{(p,q)}(M) \to \Omega_\C^{(p+1, q)}(M) \hspace{1 cm}
\barpartial:=\Pi^{(p, q+1)} \circ \diff \colon \Omega_\C^{(p,q)}(M) \to \Omega_\C^{(p, q+1)}(M)
\]
Their coordinate expressions are given by 
\begin{align*}
    \partial \alpha & = \sum_{I \bar{J}}' \derpar{\alpha_{I \bar{J}}}{z^k} \cdot \diff z^k \wedge \diff z^{i_1} \wedge \ldots \wedge \diff z^{i_p} \wedge \diff \bar{z}^{j_1} \wedge \ldots \wedge \diff \bar{z}^{j_q} \\
    \barpartial \alpha & =\sum_{I \bar{J}}' \derpar{\alpha_{I \bar{J}}}{\bar{z}^k} \cdot \diff \bar{z}^k \wedge \diff z^{i_1} \wedge \ldots \wedge \diff z^{i_p} \wedge \diff \bar{z}^{j_1} \wedge \ldots \wedge \diff \bar{z}^{j_q}
\end{align*}

The following properties also hold for \(\alpha \in \Omega_\C^k(M)\) and \(\beta \in \Omega_\C^\ell(M)\):

\begin{minipage}{0.46\textwidth}
    \begin{align}
        \partial (\alpha \wedge \beta) &= \partial \alpha \wedge \beta + (-1)^k \cdot \alpha \wedge \partial\beta. \\
        \barpartial (\alpha \wedge \beta) &= \barpartial \alpha \wedge \beta + (-1)^k \cdot \alpha \wedge \barpartial\beta. \\
        \diff \alpha &= \partial \alpha + \barpartial \alpha. \\
        \barpartial f&=0 \ \ \text{if} \ \ f \in \mathcal{O}_U \label{barparial_holomorphic_zero}
    \end{align}
\end{minipage}
\begin{minipage}{0.46\textwidth}
    \begin{align}
        \overline{\partial \alpha}&=\barpartial \bar{\alpha}. \\
        \partial \partial \alpha&=\barpartial \barpartial \alpha=0 \label{eqbarpartial2} \\
        \partial \barpartial \alpha &= -\barpartial \partial \alpha \\
        \partial f&=0 \ \ \text{if} \ \ \bar{f} \in \mathcal{O}_U
    \end{align}
\end{minipage}

% \begin{multicols}{2}
%     \begin{enumerate}
%         \item \(\partial (\alpha \wedge \beta) = \partial \alpha \wedge \beta + (-1)^k \cdot \alpha \wedge \partial\beta\).
%         \item \(\barpartial (\alpha \wedge \beta) = \barpartial \alpha \wedge \beta + (-1)^k \cdot \alpha \wedge \barpartial\beta\).
%         \item \(\diff \alpha = \partial \alpha + \barpartial \alpha\).
%         \item \(\overline{\partial \alpha}=\barpartial \bar{\alpha}\).
%         \item \(\partial \partial \alpha=\barpartial \barpartial \alpha=0\).
%         \item \(\partial \barpartial \alpha = -\barpartial \partial \alpha\).
%     \end{enumerate}
% \end{multicols}

% Condition \eqref{eqbarpartial2} means that, for each \(0 \leq p\leq n\), we have a \textit{Dolbeault complex}
% \[
% \corch{0} \xrightarrow[]{} \Omega_\C^{(p,0)}(M) \xrightarrow[]{\barpartial} \Omega_\C^{(p,1)}(M) \xrightarrow[]{\barpartial} \Omega_\C^{(p,2)}(M) \xrightarrow[]{\barpartial} \cdots \xrightarrow[]{\barpartial} \Omega_\C^{(p,n)}(M) \to \corch{0}.
% \]
% Thus, we can define the \textit{Dolbeault cohomology groups}:
% \begin{equation}
%     H^{(p,q)}(M):=\frac{\ker \parent{\barpartial\colon \Omega_\C^{(p,q)}(M) \to \Omega_\C^{(p,q+1)}(M)}}{\mathrm{Im} \parent{\barpartial \colon \Omega_\C^{(p,q-1)}(M) \to \Omega_\C^{(p,q)}(M)}}.
% \end{equation}

The exterior differential produces de Rham cohomology groups
\[
\corch{0} \xrightarrow{} \Omega_\C^0(M) \xrightarrow{\diff} \Omega_\C^1(M) \xrightarrow{\diff} \Omega_\C^2(M) \xrightarrow{\diff}\cdots \xrightarrow{\diff} \Omega_\C^n(M) \to \corch{0}, 
\]
which are called the \textit{de Rham cohomology groups with complex coefficients}\index{De Rham cohomology with complex coefficients}, and denoted \(H^k_{\mathrm{dR}}(M;\C)\). The real de Rham cohomology groups are denoted \(H^k_{\mathrm{dR}}(M;\R)\). They can be seen as real subspaces of the complex de Rham cohomology groups (classes of forms invariant under conjugation).

\section{Bundle-Valued forms} \label{section_bundle_valued_p_forms}
Let \(M\) be a smooth manifold, and let \(E\to M\) be a smooth real vector bundle. A \textit{real \(E\)-valued \(p\)-form} is a section of the real vector bundle \(\bigwedge^p T^*M \otimes E\). We define 
\[
\Omega^k_\R(U;E):=\Gamma\parent{U;{\bigwedge}^pT^*M\otimes E}
\] 
the space of all \(E\)-valued \(p\)-forms defined on an open set \(U \subset M\). Similarly, if \(E\) is a smooth complex vector bundle (observe that we do not require \(E\) to be holomorphic), we define a \textit{complex \(E\)-valued \(k\)-form} as a section of the bundle \(\bigwedge^k_\C M \otimes E\), and a \textit{complex \(E\)-valued \((p,q)\)-form} as a section of the bundle \(\bigwedge_\C^{(p,q)}M \otimes E\). We define 
\[
\Omega^k_\C(U;E):=\Gamma \parent{U; {\bigwedge}^k_\C M \otimes E} \hspace{1 cm} \text{and} \hspace{1 cm}
\Omega^{(p,q)}_\C(U;E):=\Gamma \parent{U;{\bigwedge}^{(p,q)}_\C M \otimes E}
\]
for an open set \(U \subset M\). If the open set is all of \(M\), and it is clear from the context, we shall omit it. Observe an element of \(\sigma \in \Omega_\C^{k}(U;E)\) is a sum of elements \(\omega \otimes \sigma\), where \(\omega \in \Omega_\C^{k}(M)\) and \(\sigma \in \Gamma(E)\). The Dolbeault direct sum decomposition of the space of complex \(k\)-forms implies that \(\omega\) can be uniquely decomposed as a sum of \((p,q)\)-forms satisfying \(k=p+q\). This then implies that any \(E\)-valued \(k\)-form can be uniquely decomposed as a sum of \(E\)-valued \((p,q)\)-forms. In conclusion, we have the following Whitney sum decomposition
\begin{equation} \label{eq_decomp_E_forms}
\Omega^k_\C(U; E)=\bigoplus_{p+q=k} \Omega_\C^{(p,q)}(U;E).
\end{equation}
By convention, we set \(\Omega_\C^0(M;E):=\Gamma(M;E)\). \\

There are several wedge product constructions on bundle-valued differential forms, and they are often used without any particular comment (specially in classical references). We devote this part of the section to the definition of those wedge products. Firstly, consider \(\alpha \in \Omega_\C^p(M)\) and \(\beta\in \Omega_\C^q(M; E)\) for a smooth complex vector bundle \(E\). Suppose that \(\beta=\tilde{\beta} \otimes s\) for a \(q\)-form \(\tilde{\beta} \in \Omega_\C^q(M)\) and a section \(s \in \Gamma(M;E)\) (note that every \(E\)-valued \(q\)-form is a linear combination of this kind of terms). We define 
\[
\alpha \wedge (\tilde{\beta} \otimes s):=(\alpha \wedge \tilde{\beta}) \otimes s \in \Omega_\C^{p+q}(M;E),
\]
where the \(\wedge\) symbol on the right-hand side is the usual exterior product of differential forms, and we extend it linearly to define a wedge product \(\wedge \colon \Omega_\C^p(M) \times \Omega_\C^q(M) \to \Omega_\C^{p+q}(M;E)\).

\begin{obs} \label{obs_wedge_bundle}
    To be somewhat more rigorous, both the wedge and tensor products are actually defined pointwise. Let \(m \in M\) and fix \(\alpha \in \bigwedge^p (T_m^*M)_\C\). The map 
    \[
    \begin{array}{cccc}
      \wedge_\alpha\colon &   \parent{\bigwedge^q (T_m^*M)_\C} \otimes_\C E_m & \longrightarrow & \parent{\bigwedge^{p+q} (T_m^*M)_\C} \otimes _\C E_m \\
       & (\beta, s) & \mapsto & (\alpha \wedge \beta) \otimes s 
    \end{array}
    \]
    is bilinear. Thus, by the universal property of the tensor product, it extends to a uniquely defined linear map making the following diagram commutative
    \begin{center}
    \begin{tikzcd}[row sep=normal, column sep=normal,  ampersand replacement=\&]
           {\bigwedge}^q (T_m^*M)_\C \times E_m \arrow[r, "{\displaystyle  \wedge_\alpha}"] \arrow[d, "{\displaystyle \otimes}", '] \& {\bigwedge}^{q+p} (T_m^*M)_\C \otimes_\C E_m \\
            {\bigwedge}^q(T_m^*M)_\C \otimes_\C E_m \arrow[ur, "{\displaystyle \tilde{\wedge}_\alpha}", '] \& 
        \end{tikzcd}
    \end{center}
    Then, for a general element \(\beta\in \parent{\bigwedge^q (T^*_mM_)\C} \otimes_\C E_m\), \(\alpha \wedge \beta\) should be understood as \(\tilde{\wedge}_\alpha (\beta)\). We will simply write \(\tilde{\wedge}_\alpha(\beta)\) as \(\alpha \wedge \beta\).
\end{obs}

When \(\beta \in \Omega_\C^0(M)\) we define \(\alpha \wedge \beta:=\alpha \otimes \beta\), and if \(\alpha \in \Omega_\C^0(M)=\mathscr{C}^\infty(M;\C)\), we define \(\alpha \wedge \beta:=\alpha \cdot \beta\). \\

We may also define a wedge product between \(E\) and \(E^*\)-valued forms by linear extension (a rigorous treatment is done as in remark \ref{obs_wedge_bundle}) of
\[
(\alpha \otimes \varphi) \wedge (\beta \otimes s):=(\alpha \wedge \beta) \otimes (\varphi(s))=\varphi(s) \cdot \alpha \wedge \beta \in \Omega_\C^{p+q}(M)
\]
if \(\alpha \otimes \varphi \in \Omega_\C^p(M;E^*)\) and \(\beta \otimes s \in \Omega_\C^q(M;E)\). Observe that \(\varphi(s) \in \mathscr{C}^\infty(M;\C)\), so the above product makes sense. This defines a wedge product \(\wedge\colon \Omega_\C^p(M;E^*) \times \Omega_\C^q(M;E) \to \Omega_\C^{p+q}(M)\). Similarly, if \(\alpha \otimes A \in \Omega_\C^p(M; \mathrm{End}(E))\) and \(\beta \otimes B \in \Omega_\C^q(M; \mathrm{End}(E))\), we can define a wedge product \(\wedge \colon\Omega_\C^p(M; \mathrm{End}(E)) \times \Omega_\C^q(M; \mathrm{End}(E)) \to \Omega_\C^{p+q}(M; \mathrm{End}(E))\) by linear extension of
\[
(\alpha \otimes A) \wedge (\beta \otimes B):=(\alpha \wedge \beta) \otimes (A \circ B).
\]
Finally, \(\wedge \colon \Omega_\C^p(M; \mathrm{End}(E)) \times \Omega_\C^q(M;E) \to \Omega_\C^{p+q}(M;E)\) can be defined by linear extension of
\[
(\alpha \otimes A) \wedge (\beta \otimes s):=(\alpha \wedge \beta) \otimes (A(s)).
\]
For more details, see \parencite[Ch. 4]{lee_introCM}. \\

When the bundle \(E\) is holomorphic, then there is a very special operator, called the \textit{Cauchy-Riemann operator}.
\Proposition{\textbf{Cauchy-Riemann operator}}{Let \(M\) be a complex manifold and \(\pi\colon E \to M\) be a holomorphic vector bundle. Then, there are unique operators \(\barpartial_E\colon \Omega_\C^{(p,q)}(M;E) \to \Omega_\C^{(p,q+1)}(M;E)\) satisfying the following properties: \label{Cauchy_Riemann_operator}
\begin{enumerate}
    \item If \(s \in \Gamma(M;E)\), then \(s\) is holomorphic if and only if \(\barpartial_E s=0\).
    \item If \(\alpha \in \Omega_\C^{(p,q)}(M)\) and \(\beta \in \Omega_\C^{(p',q')}(M;E)\), then
    \begin{equation} \label{prop_operador_cauchy_riemann}
    \barpartial_E(\alpha \wedge \beta):=\barpartial \alpha \wedge \beta + (-1)^{p+q} \alpha \wedge \barpartial_E \beta.
    \end{equation}
    \item \(\barpartial_E \circ \barpartial_E=0\).
\end{enumerate} 
}
\begin{proof}
    As usual for this kind of constructions, the key idea is to define the operator in terms of a local frame, and then show that the definition is independent of the chosen frame. Thus, let \((s_1, \ldots s_r)\) be a local holomorphic frame. Let \(\sigma \in \Omega_\C^{(p,q)}(M;E)\). Locally, we may write it as \(\sigma=\alpha^j \otimes s_j\) for some \(\alpha^j \in \Omega_\C^{(p,q)}(M)\). If \(\barpartial_E\) existed and satisfied conditions 2 and 3, we would have 
    \begin{equation} \label{def_barpartial_E}
    \barpartial_E \sigma = \barpartial \alpha^j \otimes s_j.
    \end{equation}
    Therefore, we define \(\barpartial_E\) in the domain of th holomorphic frame we chose by equation \eqref{def_barpartial_E}. To check that this defines \(\barpartial_E\) globally on \(M\), choose another local frame \((\tilde{s}_1, \ldots \tilde{s}_r)\) whose domain intersect that of the original frame we had. With respet to that frame, we may write \(\sigma=\tilde{\alpha}^k \otimes \tilde{s}_k\). By proposition \ref{prop_change_of_coordintes_frames}, there are holomorphic functions \(\tau^j_k\) such that \(\tilde{s}_k=\tau^j_k \cdot s_j\). Then, \(\tau^j_k\cdot \tilde{\alpha}^k=\alpha^j\). 
    \begin{align*}
        \barpartial_E\sigma & = \barpartial \alpha^j \otimes s_j = \barpartial(\tau^j_k\cdot \tilde{\alpha}^k) \otimes s_j = \parent{\barpartial (\tau^j_k) \wedge \tilde{\alpha}^k + \tau^j_k \cdot \barpartial \tilde{\alpha}^k} \otimes s_j \overset{{\color{red}(1)}}{=} \barpartial \tilde{\alpha}^k \otimes \parent{\tau^j_k \cdot s_j} \\
        & = \ \barpartial \tilde{\alpha}^k \otimes \tilde{s}_k,
    \end{align*}
    hence \(\barpartial_E\) is well-defined. Observe that in \({\color{red}(1)}\) we have used that \(\tau^j_k\) is holomorphic, so \(\barpartial \tau^j_k=0\) (see \eqref{barparial_holomorphic_zero}). Now, 3 follows immediately from the fact that \(\barpartial^2=0\). To show 1, let \(\sigma=f^j \cdot s_j\) for some functions \(f^j \in \mathscr{C}^\infty(M;\C)\). We have \(\barpartial_E \sigma= \barpartial f^j \otimes s_j\), which is \(0\) if and only if \(\barpartial f^j=0\) for each \(1 \leq j \leq r\) again by \eqref{barparial_holomorphic_zero}. Finally, to check 2, consider \(\beta \in \Omega_\C^{(p',q')}(M)\) and \(\alpha \in \Omega_\C^{(p,q)}(M)\). With respect to some local holomorphic frame, we may write \(\beta=\beta^j \otimes s_j\), \(\beta^j \in \Omega_\C^{(p',q')}(M)\). Then,
    \begin{align*}
        \barpartial_E (\alpha \wedge \beta)&=\barpartial_E ((\alpha \wedge \beta^j) \otimes s_j)=\barpartial(\alpha \wedge \beta^j) \otimes s_j=\left[ \ \barpartial \alpha \wedge \beta^j + (-1)^{p+q} \cdot \alpha \wedge \barpartial \beta^j \right] \otimes s_j \\
        & = \barpartial \alpha \wedge (\beta^j \otimes s_j) + (-1)^{p+q} \cdot \alpha \wedge \left[ \ \barpartial \beta^j \otimes s_j \right] \\
        & = \partial \alpha \wedge \beta + (-1)^{p+q} \cdot \alpha \wedge \barpartial_E \beta.
    \end{align*}
\end{proof}
We will get back to this operator.

\section{Kähler manifolds}
Let \(M\) be a complex manifold, and recall that its canonical almost complex structure \(J\) induces a complex vector bundle structure on \(T_JM\). A hermitian fiber metric in \(T_JM\) that case is a map \(h\colon T_JM \times T_JM \to \mathscr{C}^\infty(M;\C)\) satisfying 
\begin{multicols}{2}
    \begin{enumerate}
        \item \(h\) is \(\mathscr{C}^\infty(M;\R)\)-bilinear.
        \item \(h(JX, Y)=\ii h(X,Y)\). 
        \item \(h(X,JY)=-\ii h(X,Y)\).
        \item \(h(Y,X)=\overline{h(X,Y)}\).
        \item \(h(X,X)>0\) for \(X\neq 0\).
    \end{enumerate}
\end{multicols}
where \(X,Y \in \Gamma(M;T_JM)\). A more general definition is given in chapter 3 (definition \ref{def_hermitian_fiber_metric}) Observe that every hermitian fiber metric on \(T_JM\) defines a riemannian metric \(g:=\mathrm{Re}(h)\). Something slightly less obvious is the fact that \(\omega=-\mathrm{Im}(h)\) is a real form of type \((1,1)\), that is  \(\omega \in \Omega_\R^2(M) \cap \Omega_\C^{(1,1)}(M)\) \parencite[Lemma 8.2]{lee_introCM}.

% \Lemma{}{Let \(M\) be a complex manifold, and let \(h\) be a hermitian fiber metric on \(T_JM\). Then, \(\omega=-\mathrm{Im}(h)\) is a real two-form of type \((1,1)\). i.e \(\omega \in \Omega_\R^2(M) \cap \Omega_\C^{(1,1)}(M)\).}
% \begin{proof}
%     The fact that \(\omega\) is smooth and \(\mathscr{C}^\infty(M;\R)\)-bilinear follows immediately for the definition of \(h\). It is antisymmetric because 
%     \[
%     \omega(X,Y)=-\frac{1}{2 \ii}\parent{h(X,Y)-\overline{h(X,Y)}}=-\frac{1}{2 \ii} \parent{h(X,Y)-h(Y,X)},
%     \]
%     and the right-hand side of this expression is clearly antisymmetric. If we stop here, we get \(\omega \in \Omega_\R^2(M)\). We extend it by complex bilineariry to check that it is a \((1,1)\)-form. By \parencite[Problem. 4.1]{lee_introCM}, it is enough to check that \(\omega(JX, JY)=\omega(X,Y)\), which holds:
%     \[
%     \omega(JX,JY)=-\frac{1}{2 \ii} \parent{h(JX,JY)-h(JY,JX)}=-\frac{\ii\cdot(-\ii)}{2 \ii}\parent{h(X,Y)-h(Y,X)}=\omega(X,Y).
%     \]
% \end{proof}

A \textit{hermitian metric}\index{Hermitian!metric} on a complex manifold \(M\) is a riemannian metric \(g\) for which \(J\) is orthogonal, i.e. \(g(JX, JY)=g(X,Y)\). If \(g\) is a hermitian metric, we say that \(M\) is a \textit{Hermitian manifold}\index{Hermitian!manifold}. If that is the case, then, we can \textit{define} \(\omega(\cdot, \cdot):=g(J\cdot, \cdot)\), and \(h:=g-\ii \omega\) and it turns out that \(h\) is a hermitian fiber metric as defined above. Under these assumptions, we call \(\omega\) the \textit{fundamental 2-form}\index{Fundamental!2-form} of \(h\). The forms \(\omega\) and \(g\) are uniquely determined by each other. Indeed \(\omega(\cdot, \cdot)=g(J \cdot, \cdot)\) and \(g(\cdot, \cdot)=\omega(\cdot, J \cdot)\) (see \parencite[lemma 8.3 and lemma 8.7]{lee_introCM}). It also follows that every complex manifold admits a hermitian metric (see \parencite[lemma 8.5]{lee_introCM}). \\

A \((1,1)\)-form \(\omega\) on a complex manifold \(M\) is said to be \textit{positive}\index{Positive!\((1,1)\)-form} if \(\omega_p(v,J_pv)>0\) for every \(v \in T_pM \setminus \corch{0}\) and \(p \in M\). Given a complex manifold \(M\) endowed with a two-form \(\omega\), if we define \(g(\cdot, \cdot):=\omega(\cdot, J \cdot)\), then \(g\) is a hermitian form whose fundamental form is \(\omega\) if and only if \(\omega\) is a positive \((1,1)\)-form (see \parencite[Prop. 8.9]{lee_introCM})

\Definition{\textbf{Kähler manifold}}{Let \(M\) be a complex manifold, and \(g\) be a hermitian metric on \(M\). We say that \(g\) is a \textit{Kähler metric}\index{Kähler!metric} if its fundamental two-form \(\omega\) is closed. In that case, we say that \(M\) is a Kähler manifold.}

Observe that every Kähler manifold is a symplectic manifold, the fundamental form being its symplectic form. Indeed, \(\omega_p\) defines an injective map \(\omega_p^\flat\colon T_pM \to T_p^*M, v \mapsto \omega(v, \cdot)\) because if \(\omega_p^\flat(v)=0\), then \(\omega_p(v,Jv)=0\), so \(v=0\)

\begin{example}
    \(\C^n\) endowed with its usual symplectic structure (when seen as \(\R^{2n}\)) is an example of a Kähler manifold. Other examples are \(\CP^n\) and all Riemann surfaces (complex manifolds of complex dimension equal to 1). Indeed, one simply fixes a hermitian metric \(g\) (this can be done by \parencite[Lemma 8.5]{lee_introCM}), and then defines \(\omega(\cdot, \cdot)=g(\cdot, J \cdot)\). This is a positive form of type \((1,1)\) \parencite[Prop. 8.9]{lee_introCM} and it is closed because every two-form is closed on a Riemann surface (its dimension as a smooth manifold is \(2\)).
\end{example}

\begin{obs}
    In order to work with coordinates more efficiently, it is quite common to introduce the following convention. Let \(M\) be a complex manifold, and consider holomorphic coordiantes \((z^1, \ldots z^n)\). As we have seen, this yields two frames for the complexified tangent and cotangent bundles, namely, 
    \[
    \left\lbrace \derpar{}{z^1}, \ldots \derpar{}{z^n}, \derpar{}{\bar{z}^1}, \dots \derpar{}{\bar{z}^n} \right\rbrace \hspace{1 cm} \text{and} \hspace{1 cm}
    \corch{\diff z^1, \ldots \diff z^n, \diff \bar{z}^1, \ldots \diff \bar{z}^n}. 
    \]
    If we number these vector and covector fields from \(1\) to \(2n\) as we have ordered them, then, for each \(1 \leq j \leq n\) we will interpret \(\bar{j}\) to mean \(j+n\). We define \(\partial_j:=\partial/\partial z^j\) and \(\partial_{\bar{j \ }}:=\partial/\partial \bar{z}^j\). Then, we may write \(X \in \Gamma(T_\C M)\) with respect to that frame \(X=X^j \partial_j+X^{\bar{j\ }} \partial_{\bar{j \ }}\). Here, \(X^{\bar{j}}\) is interpreted as a \(\mathscr{C}^\infty(M;\C)\) function totally independent from \(X^j\). \\
    
    The same applies to the coframe for \(T_\C^*M\). For example, if \(M\) is a Kähler manifold, let \(g\) be its riemannian metric. By abusing notation, we extend it by complex bilinearity to act on complex vector fields, but keep the same notation. Then, \(g\) is a section of the bundle \(T_\C^*M \otimes T_\C^*M\). In local holomorphic coordiantes, it can be written as
    \begin{equation} \label{g_en_coord_1}
    g=g_{jk} \cdot \diff z^j \otimes \diff z^k+g_{j \bar{k}} \cdot \diff z^j \otimes \diff \bar{z}^k+g_{\bar{j \ } k} \cdot \diff \bar{z}^j \otimes \diff z^k+g_{\bar{j \ }\bar{k}} \cdot \diff \bar{z}^j \otimes \diff \bar{z}^k,
    \end{equation}
    where \(g_{j\bar{k}}=g(\partial_j, \partial_{\bar{k}})\) and analogously for all other possible combinations. The coordinate expression actually simplifies a lot in this local frame, because \(g(X,Y)=0\) if \(X,Y \in T^{(1,0)}M\) (see \parencite[Prop. 8.4]{lee_introCM}) and the indices satisfy the following symmetry conditions \(\overline{g_{jk}}=g_{\bar{j \ }\bar{k}}=g_{\bar{k} \bar{j \ }}\) and \(\overline{g_{j \bar{k}}}=g_{\bar{j \ }k}=g_{k \bar{j \ }}\). Then, equation \eqref{g_en_coord_1} simplifies to
    \[
    g=g_{j\bar{k}}\cdot \diff z^j \otimes \diff \bar{z}^k+g_{\bar{j}k} \cdot \diff \bar{z}^j \otimes \diff z^k.
    \]
\end{obs}
% For many computations, we need a hermitian fiber metric on the whole complexified tangent bundle, not only on \(T_JM\). This is accomplisehd by defining 
% \[
% \langle X, Y \rangle:=g(X, \overline{Y}),
% \]
% observe that it is sesguiliear with respect to the second entry, and its coordinate-expression is given by \(\langle Z, W \rangle:=g_{j \bar{k}} Z^j W^{\bar{k}}\) for \(Z,W \in \Gamma(M; T^{(1,0)}M)\). 
Finally, observe that the relation \(\omega(X,Y)=g(JX,Y)\) provides a coordinate expression for \(\omega\) in terms of the coframe 
\[
\omega= \ii g_{j \bar{k}} \cdot \diff z^j \wedge \diff \bar{z}^k.
\]

\section{Hodge Theory}
\subsection{The Hodge inner product}
The aim o this section is to define a hermitian inner product on the set of differential forms on a complex manifold. This can be used to define the Lefschetz operator and its dual, which is essential to define Einstein-Hermitian connections. This can be done directly on vector spaces, but we may also give all definitions directly for manifolds. For the first definitions, there is no need to consider complex manifolds, riemannian manifolds suffice. \\

Let \((M,g)\) be an \(n\)-dimensional compact, oriented riemannian manifold. The riemannian metric can be complexified by extending it by complex bilinearity so that it acts on the complexified tangent bundle \(T_\C M\). We can use the musical isomorphism \(\sharp\colon T_\C^*M \to T_\C M\) to induce a hermitian fiber metric on \(\Omega_\C^1(M)\), which is given as follows:
\begin{equation} \label{eq_def_hodge_pointwise_inner_prod}
\langle \alpha, \beta \rangle :=g\parent{\alpha^\sharp, \overline{\beta^\sharp}}.
\end{equation}

Firtsly, we will extend this inner product to higher order forms. For smooth complex-valued functions, we simply define \(\langle f,h \rangle := f\cdot \overline{h}\). For higher degree forms, consider any local orthonormal frame \(\corch{\varepsilon^1, \ldots \varepsilon^n}\) for the hermitian fiber metric \eqref{eq_def_hodge_pointwise_inner_prod}. The hermitian inner product on \(\bigwedge_\C^k M\) is defined such that \(\corch{\varepsilon^{i_1} \wedge \ldots \wedge \varepsilon^{i_k}: i_1 < \ldots < i_k}\) is orthonormal. More precisely, there is a unique hermitian fiber product on \(\bigwedge^k_\C M\) which satisfies
\[
\langle \alpha_1 \wedge \ldots \wedge \alpha_k, \ \beta_1 \wedge \ldots \wedge \beta_k \rangle =\det \parent{(\langle \alpha_i, \beta_j \rangle)_{1 \leq i,j \leq k}},
\]
for any locally defined one-forms \parencite[Lemma 9.1]{lee_introCM}. We will write \(||\beta||^2=(\beta, \beta)\). When \(M\) is compact, this yields an inner product on \(\Omega_\C^k(M)\) given as follows:
\begin{equation} \label{eq_global_Hodge_inner_prod}
    (\alpha, \beta):=\int_M \langle \alpha, \beta \rangle \ \diff V_g, \hspace{1 cm} \alpha, \beta \in \Omega_\C^k(M).
\end{equation} 
Compactness of \(M\) ensures the convergence of the integral above. Equation \eqref{eq_global_Hodge_inner_prod} defines a hermitian product on the complex vector space of complex \(k\)-forms. This is called the \textit{(global) Hodge inner product}\index{Hodge! inner product} or \(L_2\) product. The symbol \(\diff V_g\) in equation \eqref{eq_global_Hodge_inner_prod} denotes the riemannian volume on \(M\). For a Kähler manifold, \(\diff V_g=\omega^n/n!\), and the orientation with respect to which we integrate is determined by \(\omega\) {\color{red}ref}.

\begin{obs} \label{def_prod_vect_bund}
    The previous construction can be generalized to any complex vector bundle over \(M\) endowed with a hermitian fiber metric \(h\). In that case, the space of smooth sections of the bundle can be endowed with an inner product as follows:
    \[
    (\xi, \eta):=\int_M h(\xi, \eta) \ \diff V_g.
    \]
    If we did not ask for compactness, we would restrict ourselves to compactly supported sections. We will write \(|\xi|:=(\xi, \xi)^{1/2}\).
\end{obs}

\subsection{The Hodge star operator}
Again, consider a compact, oriented, riemannian manifold and set \(n=\dim M\). The Hodge star operator is a vector bundle isomorphism taking 
\(k\)-forms to \(n-k\)-forms. Just as the Hodge inner product, the construction does not use anything particular to vector bundles: it is purely linear, and everything can be done in euclidean vector spaces with a fixed orientation. The operator is defined in the following proposition (again, compactness is not needed for the definition).

\Proposition{}{Let \((M,g)\) be an oriented riemannian \(n\)-dimensional manifold. Then, there is a unique smooth vector bundle isomorphism
\[
\star\colon {\bigwedge}^k_\C M \longrightarrow {\bigwedge}^{n-k}_\C M,
\]
called the \textit{Hodge star operator}\index{Hodge!star operator} satisfying, for every \(\alpha, \beta \in \bigwedge^k_\C M\), the following relation
\begin{equation} \label{eq_Hodge_star_operator}
    \alpha \wedge \star  \overline{\beta}=\langle \alpha, \beta \rangle \ \diff V_g.
\end{equation}
}
\begin{proof}
    First, we define the operator locally in terms of a coframe, and we then check that the definition does not depend on the choice of the coframe. Therefore, let \(\corch{\varepsilon^1, \ldots \varepsilon^n}\) be a local (real) frame of \(T^*M\), so, after extending by complex linearity, we get an orthonormal frame \(\corch{\varepsilon^{i_1} \wedge \ldots \wedge \varepsilon^{i_k}: i_1 < \ldots < i_k}\) of \({\bigwedge}^k_\C M\). Let \(I=(i_1, \ldots i_k)\), and define \(\varepsilon^I:=\varepsilon^{i_1} \wedge \ldots \wedge \varepsilon^{i_k}\). Then, define
    \begin{equation} \label{eq_def_star_coframe}
    \star \ \varepsilon^I= \mathrm{sgn}(\sigma(I,I'))\cdot \varepsilon^{I'}.
    \end{equation}
    Here, \(\sigma(I,I')\) is the permutation taking \(\corch{1, \ldots n}\) to \((i_1, \ldots i_k, i'_1, \ldots i'_{n-k})\). Let us see that this definition does not depend on the choice of local coframe. To do so, we shall see that definition \eqref{eq_def_star_coframe} satisfies \eqref{eq_Hodge_star_operator}, and that the latter determines a bundle isomorphism uniquely. Since both sides of \eqref{eq_Hodge_star_operator} are linear and conjugate linear in \(\alpha\) and \(\beta\), it suffices to check it for basis forms. Hence, let \(J\) and \(K\) be two sets of ascending indices of length \(k\), and \(\alpha:=\varepsilon^J, \beta :=\varepsilon^K\). Then, \(\overline{\beta}=\beta\) because we choose a real coframe, and \(\star \beta = \mathrm{sgn}(\sigma(K,K'))\cdot \varepsilon^{K'}\) by definition. Thus, \(\beta \wedge \star \overline{\beta}=\varepsilon^K \wedge \varepsilon^{K'}=\diff V_g\). Finally, we separate cases:
    \begin{enumerate}
        \item If \(J=K\) then \(\alpha=\beta\), \(\langle \alpha, \beta\rangle = 1\) and we just showed that \eqref{eq_Hodge_star_operator} is satisfied.
        \item If \(J \neq K\) then \(\langle \alpha, \beta \rangle = 0 \) because the frame \(\varepsilon^I\), for \(I\) increasing and on length \(k\) is orthonormal. The right-hand side vanishes as well \(\alpha \wedge \beta\) has some repeated element.
    \end{enumerate}
    Finally, to show uniqueness, suppose that \(\tilde{\star}\) satisfies \eqref{eq_Hodge_star_operator} as well, and let \(\delta=\star-\tilde{\star}\). Observe that for every \(\alpha, \beta \in \Omega_\C^q(M)\) we have \(\alpha \wedge \delta \overline{\beta}=0\). Therefore,
    \[
    |\delta \beta |^2 \cdot \diff V_g = \delta \beta \wedge \star \overline{\delta \beta}= \pm \star \overline{\delta \beta} \wedge \delta\beta=0,
    \]
    so \(|\delta \beta |^2=0\) for any \(\beta\), and since we have a hermitian fiber metric, \(\delta \equiv 0.\)
\end{proof}

Moreover, the star operator satisfies \parencite[Prop. 9.8]{lee_introCM}:
\begin{align}
        \overline{\star \  \alpha} &= \star \ \overline{\alpha}. \\
        \star \star \ \alpha & = (-1)^{k\cdot (N-k)} \alpha. \\
        \langle \star \ \alpha, \star \ \beta \rangle & = \langle \alpha, \beta \rangle.
\end{align}

By \eqref{eq_Hodge_star_operator}, the global inner product can be equivalently written as
\begin{equation} \label{eq_Hodge_inner_prod_def_2}
    (\alpha, \beta)=\int_M \alpha \wedge \star \overline{\beta}. 
\end{equation}
This can be done in a Riemannian manifold using real forms. In that case, an inner product on \(\Omega_\R^p(M)\) is given by \(\int_M \alpha \wedge \star \beta\).

\begin{obs}
    Everything we have done can be done in riemannian manifolds without requiring complexification of the cotangent bundle. In that case, the pointwise inner product is defined directly on \(T^*M\) by \(\langle\alpha, \beta\rangle:=g(\alpha^\sharp, \beta^\sharp)\) (no conjugation), the Hodge star operator satisfies \(\alpha \wedge \star \beta = \langle \alpha, \beta \rangle \cdot \diff V_g\), and the global inner product can be defined as \((\alpha, \beta)=\int_M \alpha \wedge \star \beta\).
\end{obs}

In fact, one can show that, on complex manifolds, for every \(\alpha \in \Omega_\C^{(p,q)}(M)\), we have \(\star  \alpha \in \Omega_\C^{(n-p, n-q)}(M)\) \parencite[Lemma 9.23-(b)]{lee_introCM}.  \\

For bundle-valued the \(L_2\) inner product admits a direct generalization. First, for any complex vector bundle \(E \to M\), endowed with a hermitian fiber metric \(h\), the bundles \(\bigwedge^k_\C  M \otimes E\) admit a unique hermitian fiber metric which is given by sesguilinear extension of the fiber metric whose action on decomposable pairs is as follows
\[
\langle \alpha \otimes \xi, \beta \otimes \eta \rangle_{q, E}=\langle \alpha, \beta \rangle \cdot h(\xi, \eta).
\]
Abusing notation, we shall write \(\langle \cdot, \cdot \rangle\) also for the fiber metric on \(\bigwedge^k_\C  M \otimes E\). By remark \ref{def_prod_vect_bund}, we get a global inner product as well. \\

% Given a smooth complex vector bundle \(E \to M\), induced with a hermitian fiber metric \(h\), we define the \textit{vector bundle Hodge star operator}\index{Hodge!vector bundle star operator} by 
% \begin{equation}
%     \begin{array}{cccc}
%         \overline{\star}_E \colon & {\bigwedge}^{(p,q)}M \otimes E & \longrightarrow & {\bigwedge}^{(n-p, n-q)} M \otimes E \\
%         & \alpha \otimes \xi & \mapsto & \overline{\star \alpha} \otimes  \hat{h}(\xi) 
%     \end{array}
% \end{equation}
% and extending bilinearly. Here, \(\hat{h}\) is conjugate-linear bundle isomorphism \(\hat{h}\colon E \to E^* \) given by 
% \[
% \begin{array}{cccc}
%     \hat{h}(\xi)\colon & E & \longrightarrow & \C \\
%     & v & \mapsto & h_p(\xi_p, v)
% \end{array}
% \hspace{1 cm} v \in E_p.
% \]

With these definitions in mind, given a Kähler manifold \((M, \omega)\) we define the \textit{Lefschetz operator}\index{Lefschetz operator} as the complex-linear bundle homomorphism given by 
\[
\begin{array}{cccc}
    L_\omega \colon & {\bigwedge}^{(p,q)} M & \longrightarrow & {\bigwedge}^{(n-p, n-q)} M \\
    & \alpha & \mapsto & \omega_p \wedge \alpha
\end{array}
\hspace{2 cm} \alpha \in \parent{{\bigwedge}^{(p,q)} M}_p 
\]
When it acts on differential forms (which is the way one usually deals with it), it acts pointwise. There is a unique operator, called the \textit{dual Lefschetz operator} or the \textit{contraction with the Kähler form}, denoted \(\Lambda_\omega\) or \(L_\omega^\star\), which is defined uniquely by the property 
\[
\langle L_\omega \alpha, \beta \rangle  = \langle \alpha, \Lambda_\omega\beta \rangle  \hspace{1 cm} \alpha \in \Omega_\C^q(M), \beta \in \Omega_\C^{q+2}(M).
\]
Explicitly, we have \(\Lambda_\omega = \star^{-1} \circ L_\omega \circ \star\). Finally, these operators can be extended to bundle valued forms \parencite{huybrechts2005complex}. Consider a complex vector bundle \(E \to M\), and define \(L_\omega \otimes \mathrm{id}_E\) and \(\Lambda_\omega \otimes \mathrm{id}_E\) on the spaces \(\bigwedge^{(p,q)}M\). That is, these maps are the complex-bilinear extension of 
\[
\parent{L_\omega \otimes \mathrm{id}_E}(\alpha \otimes \xi) = \parent{L_\omega \alpha} \otimes \xi \hspace{2 cm} \parent{\Lambda_\omega \otimes \mathrm{id}_E}(\beta \otimes \eta) = \parent{\Lambda_\omega \beta} \otimes \eta.
\]

\begin{obs} \label{obs_hodgestar_equal_lambda}
    For two-forms on a Riemann surface \(X\), one has \(\star \alpha = \Lambda_\omega \alpha\). This follows from the fact that \(\star \omega=1\) (see \parencite[6.2.1]{munoz2020geometry} and recall that \(\diff V_g=\omega\) for Riemann surfaces), so for \(f \in \mathscr{C}^\infty(X;\C), \alpha \in \Omega^2_\C(X)\)
    \[
    \langle L_\omega f, \alpha\rangle = \langle f \cdot \omega, \alpha\rangle = f\cdot \langle \star 1, \star \star \alpha \rangle = f\cdot \langle 1, \star\alpha \rangle  =  f\cdot \overline{\star \alpha}=\langle f, \star \alpha \rangle.
    \]
\end{obs}

Finally, a lemma we will need.

\Lemma{}{Let \(M\) be an Kähler manifold of compelx dimension \(n\) and \(\alpha \in \Omega_\C^{2}(M)\). Then, \label{lemma_wedge_omega}
\[
\Lambda_\omega \alpha \cdot \omega^n = n\cdot L_\omega^{n-1}\alpha.
\] 
}
\begin{proof}
    From \parencite[Lemma 9.39]{lee_introCM} we know that \([L_\omega, \Lambda_\omega]=(k-n)\cdot \mathrm{id}\) acting on \(k\)-forms, where \([\cdot, \cdot]\) is the commutator of operators. Inductively, one shows, \([L_\omega^p, \Lambda_\omega]\alpha=p\cdot(k-n+p-1) \cdot L_\omega ^{p-1}\alpha\) for each \(k\)-form \(\alpha\). Indeed, \(p=1\) is the lemma. For the inductive step (we do not write the \(k\)-form)
    \begin{align*}
        [L_\omega^p, \Lambda_\omega] & = L_\omega^p\Lambda_\omega-\Lambda_\omega L_\omega^p =  L_\omega^p \Lambda_\omega-\Lambda_\omega L_\omega L_\omega^{p.-1} \overset{{\color{red}(1)}}{=} L_\omega^p \Lambda_\omega-(L_\omega \Lambda_\omega-(k+2(p-1)-n)\cdot \mathrm{id})L_\omega^{p-1} \\
        & = L_\omega (L_\omega^{p-1}\Lambda_\omega-\Lambda_\omega L_\omega^{p-1})+(k+2(p-1)-n) \cdot L_\omega^{p-1} \\
        & = (p-1)\cdot (k-n+p-2) L_\omega^{p-1} + (k+2(p-1)-n) L_\omega^{p-1} \\
        & = p\cdot (k-n+p-1) L_\omega^{p-1}
    \end{align*}
    In {\color{red}\((1)\)} we used that \(\Lambda_\omega L_\omega\) now acts on a \(k+2(p-1)\)-form because it acts on \(L_\omega \alpha\). Once this formula has been proven, we apply it to a two-form and \(p=n\). Thus, \([L_\omega^n, \Lambda_\omega]\alpha=n\cdot L_\omega^{n-1}\alpha\). Hence (observe that \(\omega^n\wedge \alpha=0\) because it is a form whose degree is higher than the dimension of the manifold)
    \[
    L_\omega^n \Lambda_\omega\alpha -\Lambda_\omega L_\omega^n \alpha= \omega^n \wedge \Lambda_\omega \alpha = \Lambda_\omega \alpha \cdot \omega^n = n \cdot \omega^{n-1} \wedge  \alpha.
    \]
\end{proof}